\documentclass[UTF8,a4paper]{ctexart}
\usepackage{geometry,amsmath,amssymb,bm,multicol,array,xcolor,tikz}
\geometry{top=3.2mm,bottom=3.2mm,left=3.4mm,right=3.4mm}
\pagestyle{empty}
\setlength{\parindent}{0pt}
\setlength{\parskip}{0pt}
\setlength{\columnsep}{3.0mm}
\setlength{\columnseprule}{0.18pt}
\setlength{\tabcolsep}{1.15pt}
\renewcommand{\arraystretch}{0.82}
\definecolor{B}{RGB}{0,78,158}
\definecolor{R}{RGB}{170,25,25}
\definecolor{G}{RGB}{0,105,65}
\newcommand{\pg}[1]{\begin{center}\colorbox{B}{\color{white}\bfseries #1}\end{center}\vspace{-2pt}}
\newcommand{\lab}[1]{\textbf{\textcolor{B}{#1}}}
\newcommand{\warn}[1]{\textbf{\textcolor{R}{#1}}}
\newcommand{\ok}[1]{\textbf{\textcolor{G}{#1}}}
\newcommand{\blk}[5]{\textbf{\textcolor{B}{#1}}\par
\lab{看}: #2\par
\lab{首}: #3\par
\lab{求}: #4\par
\lab{检}: #5\par\vspace{1pt}\hrule\vspace{1.2pt}}
\newcommand{\mini}[2]{\textbf{\textcolor{B}{#1}} #2\par\vspace{0.6pt}}
\newcommand{\chain}[3]{\lab{链}: #1 $\rightarrow$ #2 $\rightarrow$ #3\par}
\newcolumntype{L}[1]{>{\raggedright\arraybackslash}p{#1}}
\newcommand{\ftitle}[1]{\vspace{1pt}\textbf{\textcolor{B}{#1}}\par\vspace{0.5pt}}
\newcommand{\pictitle}[1]{\scriptsize\textbf{\textcolor{B}{#1}}}
\begin{document}
\fontsize{5.05}{5.28}\selectfont\raggedright

% ================= PAGE 1 =================
\pg{P1 总读题系统：先定模型，再套模板}
\begin{multicols}{3}
\blk{0. 全题通用四行答案}
{任何新题；特别是看不懂时。}
{第一行写模型句：“本题取\underline{\hspace{0.5cm}}为控制体/流线/截面/区域，属于\underline{\hspace{0.55cm}}模型。”}
{第二行写主方程；第三行写未知量链：已知量$\to$中间量$\to$目标量；第四行写数值与单位。}
{结尾必须查：单位；表压/绝压；方向正负；截面速度是否平均；能量是否跨泵/激波/损失。}

\blk{1. 30 秒圈题法}
{题干长、图复杂、符号多。}
{先圈五类词：物质(水/空气/理想/粘性)；对象(管/板/喷管/圆柱/闸门)；状态(定常/无旋/层流/Ma)；目标($p,Q,F,M,t$)；特殊词(泵/背压/激波/相似)。}
{按目标反推：求$p$找高度/能量/等熵；求$Q$找压差/面积/阻塞；求$F$找压力积分/CV/$C_D$；求$M$找压比/面积比/波表；求$t$找液面连续。}
{若同题出现粘性和Bernoulli，用“机械能+损失”；出现激波就分段，不能一式到底。}

\blk{2. 先判能不能用公式}
{你准备套 Bernoulli/H--P/等熵/势流/Moody 时。}
{先写适用条件。Bernoulli：定常、不可压、沿流线，实际流加损失/泵；H--P：圆管、层流、充分发展；等熵：绝热无摩擦无激波；势流：无粘无旋。}
{条件不满足就改公式：Bernoulli 加$h_L,H_p,H_T$；等熵遇激波要“波前等熵$\to$激波$\to$波后等熵”；层流/湍流先算$Re$。}
{不要把 $0.528$ 当万能压比；它只对应空气 $\gamma=1.4$、$M=1$临界状态。}

\blk{3. 截面编号法}
{图里有水箱、管路、泵、阀门、喷管、压力表。}
{第一步在图上标 1,2,3；每个截面写 $z,p,V,A$。大水面：$V\approx0$，开口表压$p=0$。}
{同一管段 $V=Q/A$；不同管径分开算。局部损失用元件所在管段速度。泵/水轮/激波/变径处必须分段。}
{表压只适合液体机械能；气体状态方程、汽蚀、蒸汽压必须用绝压。}

\blk{4. 单位换算急查}
{任何数值题。}
{先把所有量换 SI。}
{$1\,{\rm cm^2}=10^{-4}{\rm m^2}$；$Q({\rm m^3/s})=Q({\rm m^3/h})/3600$；$1{\rm kPa}=1000{\rm Pa}$；水头$m=p/(\rho g)$，水中$1{\rm kPa}\approx0.102m$。}
{所有最终答案写单位；力 N，压强 Pa/kPa，流量 $\rm m^3/s$，质量流量 kg/s。}

\blk{5. 结果异常回查}
{算出负压力、负损失、流量不守恒、Ma 支路怪、方向反。}
{先不重算，按五项查：单位、模型、边界、支路、符号。}
{负绝压：表压/绝压混；损失为负：上下游或速度头写反；质量流量不等：密度或面积错；Ma 两支：看亚声/超声支。}
{力题最后问“流体对管壁”还是“管壁对流体”；两者相反。}

\blk{6. 概念题三句模板}
{题问“简述/说明/为什么/定义/意义”。}
{第一句：按题干原词定义。}
{第二句：写物理原因或适用条件。第三句：写公式、判据或后果。}
{不写空泛话。每句都要含关键词：守恒、边界、相似、压缩、粘性、能量损失、总压变化等。}

\blk{7. 证明/推导题模板}
{题问证明、推导、由什么得到什么。}
{第一步写“在\underline{\hspace{0.45cm}}条件下，由\underline{\hspace{0.6cm}}方程出发”。}
{固定顺序：假设条件$\to$微分/积分方程$\to$边界条件/积分路径$\to$化简$\to$结论与限制。}
{最后必须写适用条件；否则推导题容易扣“理论前提”分。}

\blk{8. 查表题三行格式}
{题给 Moody、斜激波图、PM 表、阻力系数图、物性表。}
{先写“查表输入”。}
{输入：$Re,\varepsilon/d$查$\lambda$；$M_1$查正激波；$M_1,\theta$查斜激波；$\nu(M)$查PM；$Re$查$C_D$；温度查$\nu,\rho,p_v$。}
{查表后写“读得\underline{\hspace{0.5cm}}，代入\underline{\hspace{0.5cm}}”。不能只写数字。}

\blk{9. 题型入口索引}
{考场翻页时不知道找哪块。}
{先用名词定位页：液面/闸门$\to$P2；管路/泵/CV$\to$P3；粘性/边界层$\to$P4；无旋/圆柱/镜像$\to$P5；Ma/喷管/激波$\to$P6。}
{再用目标量定位块：求压强找$p$链；求流量找$Q$链；求力找$F$链；求时间找连续积分；求马赫数找压比/面积比/波表。}
{不要按“我记得哪个知识点”找；按题干实物和要求量找。}

\blk{10. 反推落笔法}
{题目只问一个结果，但不知道从哪里开始。}
{写“目标量由\underline{\hspace{0.5cm}}式给出，需要先求\underline{\hspace{0.5cm}}”。}
{例：$P=\rho gQH$ 先求$Q,H$；$F=pA$先求$p$和作用点；$\dot m=\rho VA$先求$M,T,p,\rho$；$D=C_DqA$先求$Re,C_D$。}
{先写反推链也能拿过程分；链断处写“查表/由边界条件确定”。}

\blk{11. 图形标注清单}
{有示意图、坐标、液面、管径、角度、截面。}
{第一分钟只标不算：坐标正方向、截面编号、自由面、外界大气、管径、面积、长度、角度、已知速度。}
{每个截面旁写$z,p,V,A$；每个元件旁写$\zeta$或泵/水轮；每个壁面旁写边界条件。}
{图上没标的符号不要凭感觉代；先自定义并说明。}

\blk{12. 表压/绝压分流}
{出现大气压、压力表、真空、汽化压、气体、状态方程。}
{首写“本题压力采用表压/绝压”。}
{液体能量方程两端都接大气可用表压；汽蚀、蒸汽压、气体$p=\rho RT$、等熵和激波必须用绝压。$p_{abs}=p_g+p_a$。}
{出现负表压不一定错；出现负绝压一定错。}

\blk{13. 平均速度/局部速度}
{题给$\bar u$、管截面平均、壁面附近速度、速度分布。}
{首写“$Q=A\bar u$，局部速度用$u(y)$或$u(r)$”。}
{层流圆管$u_{\max}=2\bar u$；壁律中$u(y)$不能直接当$\bar u$；动量通量严格要$\int \rho u^2dA$，常规管路用平均速度头。}
{看到上横线就当平均；没有上横线看题意是点速度还是截面速度。}

\blk{14. 先算无量纲数}
{题中出现“层流/湍流/阻力/模型/可压缩/边界层”。}
{首写判据而不是直接套公式。}
{$Re=VL/\nu$判粘性流态；$Ma=V/a$判压缩；$Fr=V/\sqrt{gL}$判重力波；$\varepsilon/d$判粗糙；$A/A^*$判喷管支路。}
{判据决定公式支路；少写判据，后面公式即使对也容易丢分。}

\blk{15. 正负号约定}
{题问方向、反力、压差、功率、力矩。}
{先在图上画正方向，并写“取向右/向上/顺时针为正”。}
{方程全按自定正方向列；算出负号就表示实际反向。压差写$\Delta p=p_1-p_2$或$p_2-p_1$必须先定义。}
{不要中途换正方向；最后用一句话说明实际方向。}

\blk{16. 没时间抢分模板}
{最后五分钟或完全不会。}
{写模型句+主方程+已知量代入位置。}
{至少给出：控制体/流线选择；守恒方程；未知链；适用条件；单位检查。数值不会算就保留表达式。}
{表达式答案必须量纲正确；不要写无来源的最终数。}

\blk{17. 常见物性默认}
{题没给物性但默认常温。}
{写“取水$\rho=1000{\rm kg/m^3}$，空气$R=287{\rm J/(kgK)},\gamma=1.4$”。}
{水动力粘度常约$10^{-3}{\rm Pa\,s}$，运动粘度约$10^{-6}{\rm m^2/s}$；重力$g=9.81{\rm m/s^2}$。}
{若题给温度或表格，优先用题给，不用默认。}

\blk{18. 最后一行答案格式}
{任何计算题。}
{最后写“故 \underline{\hspace{0.8cm}} = 数值 单位，方向/状态为\underline{\hspace{0.8cm}}”。}
{压力题写表压或绝压；力题写作用对象；喷管题写是否阻塞/是否有激波；流量题写体积流量或质量流量。}
{没有单位和状态说明，等于把最后判断分送掉。}

\end{multicols}
\ftitle{P1 全局判型表：先按题干词定位，不按记忆乱翻}
\begin{tabular}{|L{0.19\textwidth}|L{0.23\textwidth}|L{0.34\textwidth}|L{0.19\textwidth}|}\hline
\textbf{看到/图形} & \textbf{第一行} & \textbf{中间顺序} & \textbf{最后检查}\\\hline
求$Q$、流量、放空 & 写连续或机械能 & $A,d\to V(Q)\to Re/\lambda\to h_L(Q)\to Q$；孔口则$Q=C_dA\sqrt{2gH}$ & 体积/质量流量别混 \\\hline
求$p,\Delta p$、压力表 & 声明表压/绝压 & 静水用$p=\rho gh$；管路用机械能；气体用$p=\rho RT$或等熵 & 汽蚀/气体必须绝压 \\\hline
求力、反力、螺栓 & 画CV或受压面 & 静水力：$F,p_D$；动量力：$\sum F=\dot m(V_2-V_1)$；外绕力：$C_DqA$ & 问谁受力就取对应方向 \\\hline
求速度分布 & 写简化NS/势流 & 粘性：方程积分+边界；势流：$\phi,\psi$叠加+求导；边界层：剖面积分 & 局部速度不等于平均速度 \\\hline
求阻力/损失 & 先判内流外流 & 内流：$Re,\lambda,h_L$；外流：$Re,C_D,D$；平板：$C_f,A_{wet}$ & 面积选迎风还是湿面积 \\\hline
求时间 & 写非定常连续 & 液面下降：$A(H)dH/dt+Q(H)=0$；振荡：受力$\to$微分方程 & 积分上下限给正时间 \\\hline
求马赫数/喷管 & 写$Ma=V/a$ & 压比$\to M$；面积比$\to M$两支；背压$\to$是否阻塞/激波 & $0.528$只对应空气临界 \\\hline
求证明/简述 & 写条件开头 & 条件$\to$方程$\to$边界/积分$\to$结论；概念题写定义+原因+后果 & 必须写适用条件 \\\hline
给复杂图 & 先编号截面 & 每截面填$z,p,V,A$；每元件填$\zeta,H_p,H_T$；每壁面填边界条件 & 不自造未定义符号 \\\hline
不会算数值 & 写表达式链 & 模型句+主方程+代入位+待查表量；保留符号答案 & 量纲正确也能拿分 \\\hline
\end{tabular}
\ftitle{P1 落笔句库：不会也要把过程分写出来}
\begin{tabular}{|L{0.19\textwidth}|L{0.23\textwidth}|L{0.34\textwidth}|L{0.19\textwidth}|}\hline
不确定模型 & “先按\underline{\hspace{0.4cm}}模型试算” & 写适用条件；若条件不符，补损失/激波/粘性修正 & 不把修正项漏掉 \\\hline
题有“忽略” & “忽略\underline{\hspace{0.4cm}}，故……” & 删去对应项：粘性删$h_L$；高度相同删$z$；大水面删$V$ & 只删题明确说的量 \\\hline
题有“证明” & “由\underline{\hspace{0.5cm}}守恒出发” & 微分式$\to$积分式$\to$边界条件$\to$结论 & 条件写在开头 \\\hline
题有“判断” & “判据为\underline{\hspace{0.5cm}}” & 先算无量纲数或压力比；再落入支路；最后说明状态 & 不只写是/否 \\\hline
题有“最大/临界” & “临界条件是\underline{\hspace{0.5cm}}” & 管流多为$Re$或允许压差；喷管为$M=1$；边界层为$Re_{cr}$ & 临界不是任意取极值 \\\hline
题有“方向” & “取\underline{\hspace{0.4cm}}为正” & 方程列完再由结果正负判实际方向 & 不能先猜方向再改式子 \\\hline
题有“单位长度” & “以下按单位长度计算” & 面积改成长度元，如圆柱$dA=a\,d\theta$；力为$N/m$ & 最后写$N/m$ \\\hline
题有“每单位质量” & “除以$\rho$得到单位质量量” & 压强项$p/\rho$，体力为$f$，能量为$J/kg$ & 不再乘体积 \\\hline
图有对称 & “利用对称性，某分量为0” & 对称面法向速度/剪力/合力某分量可消 & 对称只消对应分量 \\\hline
题给经验式 & “按题给经验式代入” & 先确认变量定义；再代数求未知；最后验范围 & 不混用别的经验式 \\\hline
\end{tabular}
\ftitle{P1 考场纠错表：看到异常先查这里}
\begin{tabular}{|L{0.22\textwidth}|L{0.25\textwidth}|L{0.33\textwidth}|L{0.15\textwidth}|}\hline
负流量 & 先别改绝对值 & 检查正方向、压差定义、积分上下限 & 说明反向 \\\hline
压力极大/极小离谱 & 查表压绝压 & 液体能量可表压；气体/汽蚀必须绝压 & 绝压$>0$ \\\hline
损失为负 & 查上下游顺序 & 损失永远耗能，应在机械能方程右侧为正 & 重列能量式 \\\hline
速度不守恒 & 查面积/密度 & 不可压$Q$守恒；可压$\dot m$守恒 & 同截面状态 \\\hline
结果缺方向 & 补一句实际方向 & “负号说明方向与假定相反” & 不擦负号 \\\hline
查表无输入 & 先算判据 & Moody要$Re,\varepsilon/d$；波表要$M,\theta$；阻力要$Re$ & 输入写出来 \\\hline
公式太长写不下 & 写公式链 & 主方程+关键中间量+最终表达式 & 保持量纲 \\\hline
\end{tabular}
\ftitle{P1 题图触发器：看到图先定位模型}
\begin{tabular}{cccccc}
\begin{tikzpicture}[scale=.38,baseline=-1ex]\draw (0,0) rectangle (1.3,.9);\draw[->] (1.3,.45)--(2,.45);\draw (2,.2)--(2,.7);\node at (1,-.35){\pictitle{水箱出口}};\end{tikzpicture}&
\begin{tikzpicture}[scale=.38,baseline=-1ex]\draw (0,.3)--(.8,.3);\draw (.95,.3) circle(.18);\draw (1.1,.3)--(1.8,.3);\draw[->] (.1,.5)--(1.7,.5);\node at (.9,-.25){\pictitle{泵/管路}};\end{tikzpicture}&
\begin{tikzpicture}[scale=.38,baseline=-1ex]\draw[thick] (0,0)--(.9,0) arc(-90:0:.45)--(1.35,.9);\draw[->] (-.1,.2)--(.6,.2);\draw[->] (1.2,.7)--(1.8,.7);\node at (.9,-.35){\pictitle{弯管CV}};\end{tikzpicture}&
\begin{tikzpicture}[scale=.38,baseline=-1ex]\draw[thick] (0,0)--(2,0);\draw[->] (-.1,.35)--(.7,.35);\draw (0,0) .. controls (.6,.15) and (1.2,.45) .. (2,.75);\node at (1,-.35){\pictitle{边界层}};\end{tikzpicture}&
\begin{tikzpicture}[scale=.38,baseline=-1ex]\draw[->] (0,.3)--(.7,.3);\draw (.95,.3) circle(.35);\draw[->] (1.3,.7) arc(55:-220:.35);\node at (1,-.35){\pictitle{圆柱/环量}};\end{tikzpicture}&
\begin{tikzpicture}[scale=.38,baseline=-1ex]\draw (0,.1)..controls(.7,.45)and(1.2,.45)..(2,.1);\draw (0,1)..controls(.7,.65)and(1.2,.65)..(2,1);\draw[->] (-.1,.55)--(.6,.55);\node at (1,-.25){\pictitle{Laval/波}};\end{tikzpicture}
\end{tabular}
\ftitle{P1 最小符号字典：先认字母，再找模板}
\begin{tabular}{|L{0.11\textwidth}|L{0.20\textwidth}|L{0.11\textwidth}|L{0.20\textwidth}|L{0.11\textwidth}|L{0.20\textwidth}|}\hline
$p$ & 静压；液体能量可用表压 & $p_0,T_0$ & 总压/总温；可压缩常用绝压 & $p_b,p_e$ & 背压/出口内压；未阻塞常有$p_e=p_b$\\\hline
$\rho,\mu,\nu$ & 密度/动力粘度/运动粘度$\nu=\mu/\rho$ & $Q,\dot m$ & 体积流量/质量流量，$\dot m=\rho Q$ & $V,u,\bar u$ & 速度/局部速度/截面平均速度\\\hline
$z,h,H$ & 高程/液柱差/水头 & $h_L,H_p,H_T$ & 损失水头/泵扬程/水轮机水头 & $\lambda,\zeta$ & 沿程摩阻系数/局部阻力系数\\\hline
$Re$ & 惯性/粘性，先判层湍流 & $Ma$ & 速度/声速，先判可压缩状态 & $C_D,C_f,C_p$ & 阻力/摩擦/压力系数，先看定义面积\\\hline
$\phi,\psi$ & 势函数/流函数；叠加后求导 & $\Gamma,Q_s,M$ & 环量/源强/偶极强度 & $A,A^*,A_T$ & 面积/临界面积/喉部面积\\\hline
\end{tabular}
\ftitle{P1 按目标量快速翻页：先找目标，不按章节想}
\begin{tabular}{|L{0.14\textwidth}|L{0.18\textwidth}|L{0.44\textwidth}|L{0.19\textwidth}|}\hline
求$p,\Delta p$ & P2/P3/P6 & 静水用$\rho gh$；管路用机械能；气体/喷管用绝压等熵或激波 & 表压/绝压先写明\\\hline
求$Q,V,\dot m$ & P2/P3/P6 & 孔口用$C_dA\sqrt{2gH}$；管路用$Q\to Re\to\lambda\to h_L$；喷管看是否阻塞 & 截面平均速度与局部速度分开\\\hline
求$F,M,P$ & P2/P3/P5 & 静水受压面/压力体；CV动量；升力$L'=\rho U\Gamma$；功率$\rho gQH$或$M\omega$ & 力方向最后取反检查\\\hline
求$\delta,\tau,D$ & P4 & 先判内流/外流和层流/湍流；壁面剪切$\tau_w$，平板/物体阻力用系数 & 面积用湿面积或迎风面积\\\hline
求$Ma$或波后量 & P6 & 先写$Ma=V/a$；等熵、正激波、斜激波、膨胀波分开走 & 有激波必降总压\\\hline
\end{tabular}
\newpage

% ================= PAGE 2 =================
\pg{P2 静水压、测压、出流、应力：看到液面先用这页}
\begin{multicols}{3}
\blk{1. 静水压基本题}
{液体静止、闸门、水深、压强分布、压力中心。}
{首写 $p=p_0+\rho gh$，深度$h$从自由液面向下量。}
{平面合力 $F=\rho gh_C A$；压力中心 $y_D=y_C+I_C/(y_CA)$；曲面分解 $F_x=\rho gh_C A_x$，$F_y=\rho gV_p$。}
{压力方向垂直受压面；压力中心比形心深；开口液面表压为0，不是绝压为0。}

\blk{2. U 管/压差计}
{图有 U 形管、水银、油水多液体、两点压差。}
{从一个已知压强点出发，沿液柱“向下加$\rho gh$，向上减$\rho gh$”。}
{同一静止连通液体同高等压；测同一管流两点常用 $\Delta p=(\rho_m-\rho)g\Delta h$。}
{别把水银高度直接当水头；换算到目标流体压差。}

\blk{3. 小孔/水箱出流}
{大水箱、小孔、孔口、液面下降、求流量或时间。}
{首写孔口 Bernoulli：$V=C_v\sqrt{2gH}$，$Q=C_dA_o\sqrt{2gH}$。}
{定水头直接代；变水头接连续：$-A_t\,dH/dt=Q(H)$，再积分。}
{水箱大才取液面$V\approx0$；有收缩/流速系数按题给。}

\blk{4. 闸门/铰链/拉杆}
{图有矩形闸门、圆弧闸门、铰链、拉杆、求力。}
{先求水压力合力和作用点，不要先列铰链反力。}
{对铰链取矩求拉杆力；再用$x,y$力平衡求铰链反力。圆弧面用水平分力+竖直压力体。}
{矩臂要垂直距离；压力方向垂直面；表压大气部分一般相消。}

\blk{5. 应力张量题}
{题给应力张量$P$和平面法向$\bm n$，求应力矢量/法向/切向/夹角。}
{首写 $\bm p_n=P\bm n$。}
{法向分量 $p_N=\bm p_n\cdot\bm n$；切向矢量 $\bm p_T=\bm p_n-p_N\bm n$；夹角 $\cos\alpha=p_N/|\bm p_n|$。}
{$\bm n$必须是单位向量；法向分量不是矩阵对角元。}

\blk{6. 速度场判定}
{题给 $u,v,w$，问连续、有旋、势函数、流函数。}
{先写连续：不可压需 $\nabla\cdot\bm V=0$。}
{旋度 $\bm\Omega=\nabla\times\bm V$；无旋才有 $\phi$。二维不可压才有 $\psi$，$u=\psi_y,\ v=-\psi_x$。}
{先判条件再积分；不要直接求$\phi,\psi$。}

\blk{7. 给速度求 $\phi,\psi$}
{题问势函数/流函数。}
{先写先验条件：$\phi$需 $v_x-u_y=0$；$\psi$需 $u_x+v_y=0$。}
{$\phi_x=u$ 先对$x$积分，再用 $\phi_y=v$ 定未知函数；$\psi_y=u$ 先对$y$积分，再用 $-\psi_x=v$ 定函数。}
{常数可省；最后用偏导反查一遍。}

\blk{8. 非定常 U 管振荡}
{等截面 U 管、水柱振荡、不计粘性。}
{取液面位移$z$，写回复力 $2\rho gAz$。}
{质量 $\rho AL$，故 $\ddot z+(2g/L)z=0$，周期 $T=2\pi\sqrt{L/(2g)}$。}
{两液面位移方向相反；$L$ 是总液柱长度，不是一边长度。}

\blk{9. 物面边界条件}
{题问不穿透、无滑移、自由面、界面。}
{先翻译边界条件。}
{固壁无粘：$\bm V\cdot\bm n=0$；有粘再加 $\bm V=\bm V_w$；自由面 $p=p_a$；两流体界面速度连续、剪应力连续。}
{边界条件决定积分常数；漏边界条件推导很难满分。}

\blk{10. 多液体分层静压}
{容器或 U 管里有油、水、汞等多层液体。}
{首写从同一连通液面逐段走压强。}
{每跨一层都用本层密度：向下$+\rho_i gh_i$，向上$-\rho_i gh_i$；到同高同液体点令压强相等。}
{密度不能平均；每段高度必须是竖直高度，不是斜长。}

\blk{11. 斜平面受压}
{斜闸门、倾斜矩形板、求合力或压力中心。}
{首写形心深度$h_C$。}
{合力$F=\rho gh_CA$；压力中心沿板距离 $s_D=s_C+I_C/(s_CA)$，其中$h=s\sin\theta$。}
{用形心“竖直深度”算合力，用沿板坐标算作用点。}

\blk{12. 压强分布图}
{题要求画压力分布或判断最大压强位置。}
{首写静水压随深度线性增大。}
{自由面表压0；底部最大$p=\rho gh$；斜板上压力图仍按竖直深度线性变化。}
{图的箭头必须垂直受压面；不是竖直向下。}

\blk{13. 浮力/浸没体}
{物体漂浮、沉没、求浮力、稳定、排水体积。}
{首写 $F_B=\rho gV_{\rm displaced}$。}
{平衡：漂浮时$F_B=G$；悬浮时$\rho_{\rm body}=\rho$；稳定看重心$G$和浮心/稳心相对位置。}
{浮力方向竖直向上；只和排开流体体积有关。}

\blk{14. 曲面压力体}
{圆弧闸门、弯曲壁面、半圆柱受水压。}
{首写曲面力分解。}
{$F_x$等于曲面竖直投影面的静水压力；$F_y$等于曲面上方压力体液重，方向按压力体在曲面上方/下方判。}
{曲面合力作用线不是简单过形心；通常先求分力再合成。}

\blk{15. 加速容器/相对静止}
{容器水平加速或旋转，液面倾斜/抛物面。}
{首写相对静止条件：$\nabla p=\rho(\bm g-\bm a)$。}
{水平加速：$dz/dx=-a/g$；旋转：$z=\omega^2r^2/(2g)+C$。}
{这是静力学在非惯性系中用等效重力；不能套普通水平液面。}

\blk{16. 变水头出流时间}
{液面随时间下降，问多久放空/到某高度。}
{首写连续方程 $A(H)(-dH/dt)=C_dA_o\sqrt{2gH}$。}
{若水箱截面积常数：$t=\frac{2A_t}{C_dA_o\sqrt{2g}}(\sqrt{H_1}-\sqrt{H_2})$。}
{积分上下限按液面从高到低；时间必须为正。}

\blk{17. 压力中心快速判断}
{想判断答案是否合理。}
{首写压力中心一定在形心下方。}
{矩形竖直板上边在自由面：$h_D=2H/3$；完全浸没矩形用公式$h_D=h_C+I_G/(h_CA)$。}
{若算出压力中心高于形心，基本是坐标或$I_G$错。}

\blk{18. 应力平衡/体力}
{题给应力分布$P(x,y,z)$，问单位质量力$\bm f$。}
{首写静力/缓慢平衡：$\nabla\cdot P+\rho\bm f=0$。}
{求$\nabla\cdot P$：第$i$个分量为$\partial p_{ix}/\partial x+\partial p_{iy}/\partial y+\partial p_{iz}/\partial z$，再除以$-\rho$。}
{注意题中$P$正负约定；若写成$\nabla\cdot P=-\rho f$就保持到底。}

\blk{19. 圆柱坐标基础}
{题给$r,\theta,z$或轴对称流。}
{先写速度分量和连续方程。}
{不可压轴对称无旋常用 $v_r,v_\theta,v_z$；二维极坐标 $v_r=\phi_r,\ v_\theta=(1/r)\phi_\theta$，$v_r=(1/r)\psi_\theta,\ v_\theta=-\psi_r$。}
{角向真实距离是$r\,d\theta$，所以公式里常有$1/r$。}

\end{multicols}
\ftitle{P2 液体静力与基础场论：看到液面、受压面、速度场就按表走}
\begin{tabular}{|L{0.19\textwidth}|L{0.23\textwidth}|L{0.34\textwidth}|L{0.19\textwidth}|}\hline
\textbf{看到/图形} & \textbf{第一行} & \textbf{中间顺序} & \textbf{最后检查}\\\hline
竖直平板受水压 & 写$h_C,A$ & $F=\rho gh_CA\to y_D=y_C+I/(y_CA)\to$对铰链取矩 & 作用点低于形心 \\\hline
倾斜闸门 & 写竖直深度$h=s\sin\theta$ & 沿板找$s_C,I_C,A$；力用$h_C$，力臂用沿板距离 & 深度不是斜长 \\\hline
圆弧/曲面闸门 & 写$F_x,F_y$分解 & 投影面求$F_x$；压力体求$F_y$；合成大小和方向 & $F_y$方向看压力体位置 \\\hline
U 管多液体 & 从已知点走压强 & 每段按本液体密度上下加减；同高同液体等压 & 高度取竖直高度 \\\hline
孔口出流 & 写$Q=C_dA\sqrt{2gH}$ & 定水头直接算；变水头列$A_t dH/dt$积分 & 系数$C_d,C_v$别漏 \\\hline
放空时间 & 写$-A(H)dH/dt=Q(H)$ & 积分$H_1\to H_2$；面积变则保留$A(H)$ & 时间正，单位秒 \\\hline
浮体/沉体 & 写$F_B=\rho gV_{\rm disp}$ & 平衡$F_B=G$；若求稳定再看重心浮心 & 浮力只竖直向上 \\\hline
应力张量给平面 & 写$\bm p=P\bm n$ & 点乘得法向；相减得切向；再算夹角 & $\bm n$先单位化 \\\hline
速度场问连续 & 写$\nabla\cdot V$ & 等于0才不可压连续；不等于0说明不满足 & 三维别漏$w_z$ \\\hline
速度场问无旋 & 写$\nabla\times V$ & 为0才求$\phi$；二维再判$\psi$ & 先判后积分 \\\hline
求流函数 & 写$u=\psi_y,v=-\psi_x$ & 先对$y$积分，再用另一个式子定函数 & 符号最容易反 \\\hline
非惯性液面 & 写$\nabla p=\rho(g-a)$ & 水平加速液面斜率$dz/dx=-a/g$；旋转液面抛物线 & 用等效重力方向 \\\hline
\end{tabular}
\ftitle{P2 几何与边界速用表}
\begin{tabular}{|L{0.19\textwidth}|L{0.23\textwidth}|L{0.34\textwidth}|L{0.19\textwidth}|}\hline
矩形竖直板 & 写$A=bH,h_C$ & 若上边在自由面：$F=\rho gbH^2/2,\ h_D=2H/3$ & 作用点距自由面 \\\hline
完全浸没矩形 & 写$h_C,A,I_G$ & $F=\rho gh_CA$；$h_D=h_C+I_G/(h_CA)$，$I_G=bH^3/12$ & $h_D>h_C$ \\\hline
圆形闸门 & 写圆心深度 & $A=\pi R^2,\ I_G=\pi R^4/4$；代压力中心公式 & 圆心不等于压力中心 \\\hline
三角形板 & 写形心深度 & 形心在底边向顶点$1/3$高处；再用$I_G$ & 先定哪边在上 \\\hline
测压管 & 写$p_g=\rho gh$ & 管内液面高度直接给测点表压水头 & 开口连大气 \\\hline
真空高度 & 写$p_{abs}=p_a-\rho gh$ & 若求表压则$p_g=-\rho gh$ & 绝压不可负 \\\hline
压力体方向不明 & 先画曲面上方水体 & 压力体真实存在则向下；虚压力体则向上，按合力平衡判 & 别只背方向 \\\hline
速度势积分 & 写$d\phi=u\,dx+v\,dy$ & 路径积分或分步积分均可；用无旋保证路径无关 & 有旋不可用 \\\hline
流函数物理意义 & 写两流线间流量 & $\Delta\psi=q$；边界流线$\psi=C$ & 常数可任选 \\\hline
应力对称性 & 写无体偶矩时$P=P^T$ & 角动量平衡使剪应力互等，如$p_{xy}=p_{yx}$ & 只适合普通连续介质 \\\hline
\end{tabular}
\ftitle{P2 考场纠错表}
\begin{tabular}{|L{0.22\textwidth}|L{0.25\textwidth}|L{0.33\textwidth}|L{0.15\textwidth}|}\hline
压强公式无效 & 看是否静止 & $p=p_0+\rho gh$只用于静止或相对静止液体 & 动水改能量 \\\hline
深度取错 & 画自由面 & 深度永远竖直向下量，不沿斜面量 & 用$h$不是$s$ \\\hline
合力方向错 & 看受压面法线 & 静水压力垂直受压面；曲面先分力 & 不竖直乱画 \\\hline
U管符号乱 & 一步一步走 & 向下加，向上减；换液体换密度 & 同高同液等压 \\\hline
压力中心离谱 & 和形心比 & 压力中心应更深；自由面矩形为$2H/3$ & 高于形心必错 \\\hline
流函数符号反 & 回代检查 & 用$u=\psi_y,v=-\psi_x$重新求速度 & 回代最稳 \\\hline
应力法向错 & 先单位化$n$ & $\bm p=P\bm n$，再点乘$\bm n$ & 不用原长度向量 \\\hline
\end{tabular}
\ftitle{P2 题图触发器：静水/出流/应力看到图怎么落笔}
\begin{tabular}{|L{0.16\textwidth}|L{0.26\textwidth}|L{0.33\textwidth}|L{0.20\textwidth}|}\hline
\begin{tikzpicture}[scale=.42,baseline=-1ex]\draw (0,0)--(0,1)--(.5,1)--(.5,0);\draw (1,0)--(1,1)--(1.5,1)--(1.5,0);\draw[dashed] (.25,.65)--(1.25,.25);\node at(.75,-.25){\pictitle{U管/压差计}};\end{tikzpicture}
& 写“沿同一连通静止液体逐段列压强” & 从开口大气压出发，向下加$\rho gh$、向上减$\rho gh$，换液体就换$\rho$ & 同高同液等压；气体题不用表压乱代\\\hline
\begin{tikzpicture}[scale=.42,baseline=-1ex]\draw (0,0) rectangle (1.2,1);\draw (.15,.65)--(1.05,.65);\draw[->] (1.2,.25)--(1.9,.25);\node at(.9,-.25){\pictitle{孔口/放水}};\end{tikzpicture}
& 写“取自由液面到孔口列Bernoulli” & $V=C_v\sqrt{2gH}$，$Q=C_dA\sqrt{2gH}$；若水位变，列$-A(H)dH/dt=Q(H)$ & $H$ 是自由面到孔口的竖直高差\\\hline
\begin{tikzpicture}[scale=.42,baseline=-1ex]\draw[thick] (0,0)--(1.2,.35);\draw[->] (.6,.2)--(.35,.9);\draw[dashed] (.6,.2)--(.9,.95);\node at(.7,-.25){\pictitle{斜闸门}};\end{tikzpicture}
& 写“平面静水合力$F=\rho gh_CA$” & 找形心深度$h_C$，再算压力中心$h_D=h_C+I_G/(h_CA)$，最后分解力 & 深度竖直量；$h_D$ 比形心更深\\\hline
\begin{tikzpicture}[scale=.42,baseline=-1ex]\draw[thick] (0,0) arc(180:360:.6); \draw[->] (.6,.2)--(.6,-.45);\draw[->] (.2,-.15)--(.9,-.15);\node at(.6,-.7){\pictitle{曲面/弧门}};\end{tikzpicture}
& 写“曲面压力分水平/竖直分量” & $F_x$ 用投影面积；$F_y$ 用压力体重量；再合成大小和方向 & 压力体方向按真实/虚拟水体判\\\hline
\begin{tikzpicture}[scale=.42,baseline=-1ex]\draw (0,0)--(1.2,.6);\draw[->] (.6,.3)--(1.05,.8);\draw[->] (.6,.3)--(.2,.75);\node at(.65,-.25){\pictitle{应力平面}};\end{tikzpicture}
& 写“单位法向量$\bm n$，应力矢量$\bm p=P\bm n$” & 先单位化$\bm n$，矩阵乘法得$\bm p$，点乘$\bm n$得法向，余量为切向 & 对称张量可用$p_{xy}=p_{yx}$\\\hline
\end{tabular}
\ftitle{P2 零基础照抄骨架}
\begin{tabular}{|L{0.16\textwidth}|L{0.25\textwidth}|L{0.36\textwidth}|L{0.18\textwidth}|}\hline
静水合力题 & “取自由液面为基准，形心深度为$h_C$” & $F=\rho gh_CA\to h_D=h_C+I_G/(h_CA)\to$ 分解/求矩 & 合力垂直受压面\\\hline
曲面受压题 & “将总压力分为水平$F_x$和竖直$F_y$” & $F_x=$投影面压力；$F_y=$压力体重量；合力$F=\sqrt{F_x^2+F_y^2}$ & 方向由压力体判\\\hline
孔口/放空题 & “自由液面到孔口列Bernoulli并用连续方程” & 定水头求$Q$；变水头写$-A(H)dH/dt=Q(H)$后积分 & 时间下限上限别反\\\hline
速度场题 & “先算$\nabla\cdot V$和$\nabla\times V$” & 连续性$\to$不可压；无旋$\to$可求$\phi$；二维再用$u=\psi_y,v=-\psi_x$求$\psi$ & 先判后积分\\\hline
应力张量题 & “$\bm p=P\bm n$，且$\bm n$必须单位化” & 法向应力$p_n=\bm p\cdot\bm n$；切向$\bm p_t=\bm p-p_n\bm n$；夹角用点乘 & 单位向量最容易漏\\\hline
\end{tabular}
\newpage

% ================= PAGE 3 =================
\pg{P3 Bernoulli、管路、泵、水轮、控制体受力}
\begin{multicols}{3}
\blk{1. 理想 Bernoulli}
{无粘、定常、不可压，沿流线；求$p,V,z$。}
{首写 $p/\rho g+V^2/2g+z=C$。}
{选两截面；大水面$V=0$；开口表压$p=0$；由能量差求速度或压强。}
{跨泵、水轮机、局部损失、长管摩擦时不能用纯 Bernoulli。}

\blk{2. 机械能方程}
{实际管路、泵、阀门、弯头、粗糙管、求$Q,p,H$。}
{首写 $p_1/\rho g+V_1^2/2g+z_1+H_p=p_2/\rho g+V_2^2/2g+z_2+h_L+H_T$。}
{$h_L=\sum\lambda L V^2/(2gd)+\sum\zeta V^2/(2g)$；$V=Q/A$；$\lambda$由$Re,\varepsilon/d$定。}
{泵加能放左边，水轮取能放右边；各损失用所在管段速度。}

\blk{3. 管路阻力流程}
{给$Q,d,L,\nu,\varepsilon,\zeta$，求压降或泵功。}
{首写 $V=Q/A,\ Re=Vd/\nu$。}
{层流 $\lambda=64/Re$；湍流查 Moody 或用近似；沿程 $h_f=\lambda(L/d)V^2/(2g)$；局部 $h_\zeta=\zeta V^2/(2g)$。}
{Moody 的$\lambda$通常是 Darcy 系数；不要和 Fanning 系数混。}

\blk{4. 已知压差求流量}
{水平管、水槽连管、给水头差求$Q$。}
{首写能量差 $\Delta H=h_L(Q)$。}
{先假设流态算$Q$；再算$Re$验证。湍流要迭代或查图；粗糙度由实际$Q$反推$\lambda$再反查$\varepsilon/d$。}
{最后检查假设流态是否自洽。}

\blk{5. 泵安装高度/汽蚀}
{吸水泵、要求最高安装高度、汽化压、安全余量。}
{从吸水池液面列到泵入口：$p_a/\rho g+z_0=p_B/\rho g+z_B+V^2/2g+h_L$。}
{限制 $p_{B,abs}\ge p_v+\Delta p_{\rm safe}$；或 $NPSH_a=(p_B-p_v)/\rho g+V^2/2g$。}
{汽蚀必须用绝压；水温给出时查$p_v,\nu$。}

\blk{6. 水轮机输出功率}
{管路中有水轮机，求输出功率。}
{首写机械能方程，把水轮机取走的水头记 $H_T$。}
{$H_T=H_{\rm up}-H_{\rm down}-h_L$；轴功率 $P=\eta\rho gQH_T$。}
{水轮机不是泵；符号反了会功率方向错。}

\blk{7. 控制体受力}
{弯管、喷嘴、射流冲板、法兰螺栓、支座反力。}
{先画 CV，标压力、重力、支座力、流入流出速度。}
{动量方程：$\sum\bm F=\sum_{\rm out}\rho Q\bm V-\sum_{\rm in}\rho Q\bm V$。分$x,y$写；压力力推入CV。}
{若题问“流体对管壁/螺栓”，最后取反。出口射大气用表压0。}

\blk{8. 喷嘴/弯管法兰}
{入口出口面积不同，求轴向力。}
{先用能量求 $p_1,p_2,V_1,V_2$。}
{再用 $x$向动量：$p_1A_1-p_2A_2+R_x=\rho Q(V_2-V_1)$。}
{压力用表压；$R_x$是外界对流体，题问法兰受力取相反。}

\blk{9. 动量矩/旋转机械}
{水轮机、泵叶轮、喷流打转轮，求转矩或功率。}
{首写角动量方程。}
{$\sum M_O=\sum_{\rm out}\rho Q(rV_\theta)-\sum_{\rm in}\rho Q(rV_\theta)$；功率 $P=M\omega$。}
{只取切向速度分量；半径和旋转方向决定正负。}

\blk{10. 串联管路}
{一条路上多段不同$d,L,\varepsilon$，同一流量。}
{首写 $Q_1=Q_2=\cdots=Q$。}
{各段分别算$V_i=Q/A_i,Re_i,\lambda_i,h_{Li}$；总损失$h_L=\sum h_{Li}$。}
{不同直径不能共用速度头；局部损失用对应位置速度。}

\blk{11. 并联管路}
{两条或多条支路连接同两节点。}
{首写 $Q=\sum Q_i$，且各支路水头损失相等。}
{对每支路写$h_{L,i}(Q_i)$；用相等水头+流量守恒求$Q_i$。}
{并联不是损失相加；相加的是流量。}

\blk{12. 突扩/突缩/阀门}
{图有突然扩大、缩小、入口、出口、阀门。}
{首写局部损失 $h_\zeta=\zeta V^2/(2g)$。}
{突扩可用 $h=(V_1-V_2)^2/(2g)$；其余按题给$\zeta$。入口/出口/弯头/阀门全部加到总损失。}
{题给$\zeta$时不要再用经验公式重复计算。}

\blk{13. 泵功率与效率}
{给泵功率、效率、流量，或求电机功率。}
{首写水力功率 $P_h=\rho gQH_p$。}
{泵给流体的水头 $H_p=P_{\rm to\,fluid}/(\rho gQ)$；若给轴功率，$P_{\rm to\,fluid}=\eta P_{\rm shaft}$；若求输入功率，$P_{\rm in}=P_h/\eta$。}
{分清“泵给水的功率”和“电机输入功率”。}

\blk{14. 文丘里/孔板流量计}
{管中收缩喉部、压差计、求流量。}
{首写连续 $A_1V_1=A_2V_2=Q$，再写 Bernoulli 加系数。}
{理想：$\Delta p=\frac12\rho(V_2^2-V_1^2)$；解得$Q$，实际乘流量系数$C_d$。}
{压差计读数先换成$\Delta p$；喉部速度大，压强小。}

\blk{15. 射流冲板}
{自由射流冲击平板、斜板、叶片。}
{首写 CV 动量方程，入口压力常为大气表压0。}
{力约等于流量乘速度变化：$\bm F_{\rm on\,fluid}=\rho Q(\bm V_{\rm out}-\bm V_{\rm in})$；板受力取反。}
{出口若沿板流走，只保留相应方向动量变化。}

\blk{16. 选截面原则}
{不知道 Bernoulli 两端取哪里。}
{选已知量最多、速度最简单、压力已知的位置。}
{常选：大水面、自由出口、压力表截面、泵前泵后、管径变化前后。跨越元件时把元件能量或损失写进去。}
{不要在弯头内部随便取截面；取直管截面更稳。}

\blk{17. 压力表读数处理}
{图上有压力表、真空表、测压管。}
{首写表压关系。}
{压力表读数$p_g=p_{abs}-p_a$；真空表读数$p_vac=p_a-p_{abs}$；测压管高度$p_g=\rho gh$。}
{液体机械能可用表压，但若后面接汽蚀/气体状态必须转绝压。}

\blk{18. 管路未知量优先级}
{管路题未知很多。}
{先求几何量：面积、速度、$Re$。}
{再求阻力系数$\lambda,\zeta$；再求损失$h_L$；最后代机械能方程求目标$p,Q,H,P$。}
{不要一上来解机械能方程；里面的损失通常依赖$Q$。}

\blk{19. 控制体质量守恒}
{多入口多出口、喷嘴混合、分流。}
{首写 $\sum \dot m_{in}=\sum \dot m_{out}$。}
{不可压同一流体：$\sum Q_{in}=\sum Q_{out}$；可压或不同密度：用$\dot m=\rho VA$。}
{可压缩题不能直接令体积流量守恒。}

\end{multicols}
\ftitle{P3 管路与动量：先选截面，再决定能量还是动量}
\begin{tabular}{|L{0.19\textwidth}|L{0.23\textwidth}|L{0.34\textwidth}|L{0.19\textwidth}|}\hline
\textbf{看到/图形} & \textbf{第一行} & \textbf{中间顺序} & \textbf{最后检查}\\\hline
水箱到出口 & 列机械能方程 & 大水面$V=0$；出口表压0；加沿程/局部损失；解$Q$或$V$ & 出口速度用出口面积 \\\hline
管路有泵 & 写$+H_p$ & $Q\to V\to Re\to\lambda\to h_L$；由能量求$H_p$；功率$P=\rho gQH_p/\eta$ & 泵是加能 \\\hline
管路有水轮机 & 写$+H_p$左、$+H_T$右 & 上游能量-下游能量-损失=$H_T$；$P=\eta\rho gQH_T$ & 水轮是取能 \\\hline
吸水泵汽蚀 & 用绝压列到泵入口 & $p_B$由能量求；限制$p_B\ge p_v+\Delta p$；反求安装高度 & 不用表压判汽蚀 \\\hline
串联多段管 & 写$Q$相同 & 分段算$V_i,Re_i,\lambda_i,h_i$；总损失相加 & 每段用自己的$d,V$ \\\hline
并联支路 & 写$\Delta H_1=\Delta H_2$ & 每支路$h_i(Q_i)$；再用$\sum Q_i=Q$求解 & 损失不相加 \\\hline
管径突然变化 & 写局部损失 & 突扩用$(V_1-V_2)^2/2g$；其余用$\zeta V^2/2g$ & $\zeta$速度基准看题给 \\\hline
文丘里/孔板 & 连续+能量 & 压差计$\to\Delta p$；$A_1V_1=A_2V_2$；解$Q$乘系数 & 喉部压强更小 \\\hline
弯管/喷嘴反力 & 画CV & 先求$p,V,Q$；再分方向列动量方程；最后取反给管壁力 & 压力用表压 \\\hline
射流冲板 & 写$\rho Q\Delta V$ & 入口速度已知；出口按板面方向；板受力为流体受力反向 & 出口大气表压0 \\\hline
旋转叶轮 & 写角动量 & 只取$V_\theta$；$\sum M=\dot m(r_2V_{\theta2}-r_1V_{\theta1})$ & 功率$P=M\omega$ \\\hline
只给压差求$Q$ & 写$\Delta H=h_L(Q)$ & 先假设流态迭代；算$Re$验证；必要时查Moody & 最后说明流态 \\\hline
\end{tabular}
\ftitle{P3 元件与能量链速用表}
\begin{tabular}{|L{0.19\textwidth}|L{0.23\textwidth}|L{0.34\textwidth}|L{0.19\textwidth}|}\hline
入口损失 & 写$h_e=\zeta_eV^2/2g$ & 尖锐入口常$\zeta\sim0.5$，题给优先 & 用管内速度 \\\hline
出口损失 & 写$h_{out}=V^2/2g$ & 若流入大水池，出口动能耗散作损失 & 不再保留出口速度头 \\\hline
弯头损失 & 写$h_b=\zeta_bV^2/2g$ & 多个弯头逐个相加 & 速度基准同管段 \\\hline
阀门损失 & 写$h_v=\zeta_vV^2/2g$ & 半开阀$\zeta$可能很大；按题给 & 不和沿程混 \\\hline
粗糙管查图 & 写输入$Re,\varepsilon/d$ & 读$\lambda$；带入沿程损失；必要时迭代 & Darcy系数 \\\hline
泵曲线交点 & 写系统曲线$H_s(Q)$ & 泵曲线$H_p(Q)$与系统曲线交点为工况 & 两边同一$Q$ \\\hline
虹吸管 & 写最高点绝压限制 & 能量求最高点$p$；要求$p_{abs}>p_v$ & 最高点最危险 \\\hline
喷嘴出口速度 & 写能量到出口 & $V_e=\sqrt{2g(\Delta H-h_L)}$；再$\dot m=\rho A_eV_e$ & 出口表压常0 \\\hline
管路输出功 & 写$P=pQ$或$\rho gQH$ & 压强功率用压差；水头功率用水头差 & 单位W \\\hline
自由液面差 & 写可用水头$\Delta z$ & 从高水面到低水面，压力均大气，速度近0 & 只剩高度差和损失 \\\hline
\end{tabular}
\ftitle{P3 考场纠错表}
\begin{tabular}{|L{0.22\textwidth}|L{0.25\textwidth}|L{0.33\textwidth}|L{0.15\textwidth}|}\hline
泵/水轮符号混 & 看能量方向 & 泵给流体加能；水轮从流体取能 & 左泵右轮 \\\hline
局部损失漏 & 圈所有元件 & 入口、出口、弯头、阀门、突扩突缩都可能有$\zeta$ & 逐项列和 \\\hline
不同管径 & 分段算速度 & 每段$V_i=Q/A_i$，损失用本段速度 & 不共用$V$ \\\hline
动量方程方向乱 & 画流入流出箭头 & $\sum F=\dot m V_{out}-\dot m V_{in}$分量写 & 最后取反 \\\hline
压力力漏面积 & 压力乘截面 & $pA$作用在CV截面；自由射流出口表压0 & 表压统一 \\\hline
功率单位错 & 水头转功率 & $P=\rho gQH$，压差功率$P=\Delta p\,Q$ & W=Pa m3/s \\\hline
迭代题卡住 & 先假设$\lambda$ & 算$Q\to Re\to\lambda$，重复一次通常够过程分 & 说明迭代 \\\hline
\end{tabular}
\ftitle{P3 题图触发器：管路题先画截面和能量方向}
\begin{tabular}{|L{0.16\textwidth}|L{0.26\textwidth}|L{0.33\textwidth}|L{0.20\textwidth}|}\hline
\begin{tikzpicture}[scale=.42,baseline=-1ex]\draw (0,.8)--(0,0)--(1.1,0);\draw (.2,.65)--(.9,.65);\draw[->] (1.1,.25)--(1.8,.25);\node at(.9,-.25){\pictitle{水箱--出口}};\end{tikzpicture}
& 写“1取大水面，2取出口，列机械能方程” & $V_1\approx0,p_1=p_2=0$表压；有损失就加$h_L$；求$V,Q$ & 出口面积用喷口/孔口面积\\\hline
\begin{tikzpicture}[scale=.42,baseline=-1ex]\draw (0,.35)--(.7,.35);\draw (.95,.35) circle(.2);\node at(.95,.35){\tiny P};\draw (1.15,.35)--(1.9,.35);\draw[->] (.15,.55)--(1.75,.55);\node at(.95,-.25){\pictitle{泵管路}};\end{tikzpicture}
& 写“能量方程左侧加泵扬程$H_p$” & 先由$Q$算$V,Re,\lambda$，再算沿程+局部损失，最后求$H_p$或功率 & 泵给流体加能，符号别反\\\hline
\begin{tikzpicture}[scale=.42,baseline=-1ex]\draw[thick] (0,0)--(.8,0) arc(-90:0:.45)--(1.25,.8);\draw[->] (-.1,.15)--(.55,.15);\draw[->] (1.05,.65)--(1.65,.65);\node at(.85,-.35){\pictitle{弯管受力}};\end{tikzpicture}
& 写“取弯管内流体为CV，列$x,y$动量方程” & 先算$pA$和$\dot mV$，分方向列$\sum F$，解壁面对流体力，再取反 & 问“管受力”要取相反力\\\hline
\begin{tikzpicture}[scale=.42,baseline=-1ex]\draw (0,.2)--(.55,.2);\draw (.75,.2) circle(.18);\draw (.95,.2)--(1.6,.2);\draw[->] (.2,.45)--(1.45,.45);\node at(.8,-.25){\pictitle{水轮机}};\end{tikzpicture}
& 写“能量方程右侧取出$H_T$” & 上游总水头减下游总水头减损失，得$H_T$；功率$P=\eta\rho gQH_T$ & 水轮机从流体取能\\\hline
\begin{tikzpicture}[scale=.42,baseline=-1ex]\draw (0,.1)--(.55,.1)--(.55,.6)--(1.35,.6)--(1.35,.1)--(1.9,.1);\draw[->] (.1,.35)--(1.8,.35);\node at(.95,-.25){\pictitle{虹吸/最高点}};\end{tikzpicture}
& 写“最高点验绝压$p_{abs}>p_v$” & 对最高点列能量求$p$；再和汽化压比较；若求高度，令$p=p_v+\Delta p$反解 & 必须用绝压，不用表压\\\hline
\end{tabular}
\ftitle{P3 零基础照抄骨架}
\begin{tabular}{|L{0.16\textwidth}|L{0.25\textwidth}|L{0.36\textwidth}|L{0.18\textwidth}|}\hline
普通管路求$Q$ & “取1、2截面，列机械能方程” & 先猜流态$\to Q/A=V\to Re\to\lambda\to h_L(Q)$，代回能量方程，必要时迭代 & 写出所有$\zeta$和$L/d$\\\hline
泵/水轮机 & “泵加$H_p$，水轮机取出$H_T$” & 能量方程求水头；泵功率$P=\rho gQH_p/\eta$；水轮机$P=\eta\rho gQH_T$ & 效率乘除别反\\\hline
吸水高度 & “入口最低允许压力用绝压约束” & 水面到泵入口列能量；令$p_{in}=p_v+\Delta p$反求最大高度 & 汽化压题绝压\\\hline
弯管/喷嘴力 & “取管内流体为CV，分$x,y$列动量” & 压力力$+pA$，质量流量$\dot m=\rho Q$，列$\sum F=\dot m(V_{out}-V_{in})$ & 管受力=流体受力相反\\\hline
并联系统 & “各支路水头损失相等，总流量相加” & 写$h_{L1}(Q_1)=h_{L2}(Q_2)$和$Q_1+Q_2=Q$；解后验$Re$ & 不能把损失相加\\\hline
\end{tabular}
\newpage

% ================= PAGE 4 =================
\pg{P4 粘性流动、壁律、边界层、外绕阻力}
\begin{multicols}{3}
\blk{1. 圆管层流 H--P}
{圆管、层流、充分发展、求速度分布/流量/压降。}
{首写 $Re<2300$ 且充分发展。}
{$u(r)=\frac{-dp/dx}{4\mu}(R^2-r^2)$；$\bar u=u_{\max}/2$；$Q=\pi R^4\Delta p/(8\mu L)$；$\lambda=64/Re$。}
{H--P 不能用于湍流和入口发展段。}

\blk{2. 平板 Couette}
{两平板间流体，上板速度$U$，无压差。}
{首写 $\mu d^2u/dy^2=0$。}
{边界 $u(0)=0,u(h)=U$；得 $u=Uy/h$，剪应力 $\tau=\mu U/h$。}
{方向：壁面对流体剪力与相对运动相反；题问板受力取反。}

\blk{3. Couette + Poiseuille}
{两板间有上板运动，又有压强梯度/重力沿流向。}
{首写 $\mu u''=dp/dx-\rho g_x$。}
{积分两次，用两个壁面无滑移定常数；流量 $Q=b\int_0^h udy$。}
{若中点速度为零，用它反求$U$或压强梯度。}

\blk{4. 两层流体 Couette}
{两种粘度流体分层，上板运动，下板固定。}
{首写每层 $u_i=A_i y+B_i$。}
{边界：两壁无滑移；界面速度连续；界面剪应力连续 $\mu_1u_1'=\mu_2u_2'$。}
{不要令两层速度梯度相等；相等的是剪应力。}

\blk{5. 湍流壁律/边界层厚度}
{题给壁面剪切、$u_\tau,y^+$、壁律、求边界层厚度。}
{首写 $u_\tau=\sqrt{\tau_w/\rho}$，$y^+=yu_\tau/\nu$，$u^+=u/u_\tau$。}
{粘性底层 $u^+=y^+$；对数层 $u^+=2.5\ln y^++5.5$；边界层外缘常取$y^+\approx30$或$u\approx U$。}
{$\bar u$ 是截面平均速度，不等于局部$u(y)$。}

\blk{6. 平板边界层阻力}
{零攻角平板、给$U,L,b,\nu$，求阻力/厚度。}
{首写 $Re_L=UL/\nu$。}
{层流：$\delta=5L/\sqrt{Re_L}$，$\bar C_f=1.328/\sqrt{Re_L}$；湍流常用 $\bar C_f=0.074/Re_L^{1/5}$。$D=\bar C_f(0.5\rho U^2)A_{wet}$。}
{单面$A=Lb$，双面$2Lb$；若给$Re_{cr}$，先算$x_{cr}=Re_{cr}\nu/U$分段。}

\blk{7. 边界层积分题}
{给速度剖面 $u/U=f(y/\delta)$，求$\delta,\theta,C_D$。}
{设 $\eta=y/\delta$。}
{$\delta^*=\delta\int_0^1(1-f)d\eta$；$\theta=\delta\int_0^1 f(1-f)d\eta$；$\tau_w=\mu U f'(0)/\delta$；零压梯度 $\tau_w=\rho U^2d\theta/dx$。}
{最后阻力：单面 $C_D=2\theta(L)/L$。}

\blk{8. 分离判断}
{问边界层分离、逆压梯度、尾流、阻力危机。}
{首写逆压梯度 $dp/dx>0$ 使近壁流体减速。}
{分离点条件 $\tau_w=\mu(\partial u/\partial y)_0=0$；之后近壁可能回流。湍流抗逆压梯度强，分离后移，阻力危机时$C_D$骤降。}
{理想势流阻力为0不能解释真实阻力；真实阻力看粘性和分离。}

\blk{9. 外绕阻力}
{球、圆柱、汽车、阻力系数图。}
{首写 $Re=UL/\nu$，查 $C_D(Re)$。}
{$D=C_D(0.5\rho U^2)A_{\rm front}$；终速题令重力差=阻力。低$Re$小球可用 Stokes：$D=3\pi\mu dU$。}
{外绕阻力面积通常是迎风面积，不是湿面积；平板摩擦才用湿面积。}

\blk{10. 圆管湍流速度分布}
{题给$u/u_{\max}=(1-r/R)^{1/n}$或$n=7$。}
{首写湍流经验速度分布。}
{常用$1/7$律：$u/u_{\max}=(1-r/R)^{1/7}$；截面平均速度与最大速度由积分或题给关系取。}
{这是湍流经验式，不可和层流抛物线混用。}

\blk{11. 水力光滑/粗糙}
{题问光滑管、粗糙管、粗糙度影响。}
{首写 $Re,\varepsilon/d$。}
{层流$\lambda$只与$Re$有关；湍流光滑区主要随$Re$；完全粗糙区主要随$\varepsilon/d$。}
{粗糙度$h$和相对粗糙度$\varepsilon/d$不能混。}

\blk{12. 入口发展长度}
{题问充分发展前后、入口段。}
{首写入口段速度剖面还在变化。}
{层流估算 $L_e/d\approx0.05Re$；湍流估算常取$L_e/d\sim 10\!-\!60$。充分发展后$u=u(r)$，$dp/dx$常数。}
{H--P 只在入口段之后用。}

\blk{13. 倾斜平板层流}
{斜板间液膜、重力沿流向、上板运动。}
{首写流向动量方程 $\mu u''+\rho g\sin\theta-dp/dx=0$。}
{按边界积分；若无压差，重力项驱动；若有板速，叠加 Couette 项。}
{角度方向要看流向，$\sin\theta$符号随坐标正向变。}

\blk{14. 油膜/窄缝阻力}
{轴承、活塞、薄油膜、窄缝剪切力。}
{首写近似平板 Couette 或 Poiseuille。}
{纯剪切 $F=\tau A=\mu U A/h$；有压差再加抛物线流量项。}
{间隙$h$很小，单位常是 mm，要先换 m。}

\blk{15. 小球沉降/悬浮}
{小球在流体中下落、终速、$Re<1$。}
{首写终速平衡：重力-浮力=阻力。}
{低$Re$用 Stokes：$3\pi\mu dU=(\rho_s-\rho)g(\pi d^3/6)$，得$U=(\rho_s-\rho)gd^2/(18\mu)$。}
{先验算$Re=Ud/\nu<1$；否则要查$C_D$。}

\blk{16. 雷诺平均/湍流应力}
{题提脉动、时均、雷诺应力。}
{首写 $u=\bar u+u'$，且$\overline{u'}=0$。}
{雷诺应力项来自脉动动量通量，如$-\rho\overline{u'v'}$；总剪应力=粘性剪应力+湍流剪应力。}
{不是新流体性质，是平均方程多出的动量输运项。}

\blk{17. 压强梯度与外流速度}
{边界层外有$U(x)$，问顺压/逆压。}
{首写外流 Bernoulli：$dp/dx=-\rho U\,dU/dx$。}
{若$U$增大，则$dp/dx<0$顺压，边界层稳定；若$U$减小，则$dp/dx>0$逆压，易分离。}
{分离不是因为速度大，而是近壁动量不足以抗逆压。}

\blk{18. 平板层/湍混合}
{给$Re_{cr}$，从前缘先层流后湍流。}
{首写 $x_{cr}=Re_{cr}\nu/U$。}
{若$L<x_{cr}$全层流；若$L>x_{cr}$，前段层流、后段湍流，阻力分段相加或用带修正的平均$C_f$。}
{不要把全板都按湍流或层流算。}

\blk{19. 边界层题落笔}
{题给速度剖面、要求厚度/阻力。}
{首写三厚度定义：$\delta,\delta^*,\theta$。}
{先代剖面积分得$\delta^*,\theta$；再用动量积分方程求$\delta(x)$；最后求$\tau_w,D$。}
{积分变量换成$\eta=y/\delta$可避免漏$\delta$。}

\end{multicols}
\ftitle{P4 粘性与边界层：先判流态，再选速度模型}
\begin{tabular}{|L{0.19\textwidth}|L{0.23\textwidth}|L{0.34\textwidth}|L{0.19\textwidth}|}\hline
\textbf{看到/图形} & \textbf{第一行} & \textbf{中间顺序} & \textbf{最后检查}\\\hline
圆管层流 & 写$Re<2300$ & $u(r)$抛物线；$\bar u=u_{\max}/2$；$Q=\pi R^4\Delta p/(8\mu L)$ & 充分发展才用 \\\hline
圆管湍流 & 写$Re,\varepsilon/d$ & 查$\lambda$或用经验式；$h_f=\lambda L V^2/(2gd)$ & 不能用H--P \\\hline
两平板剪切 & 写$\mu u''=0$ & 两边界定线性速度；$\tau=\mu U/h$；力$F=\tau A$ & 板受力方向取反 \\\hline
压差驱动薄层 & 写$\mu u''=dp/dx$ & 积分两次；无滑移定常数；流量积分 & $dp/dx$符号跟坐标一致 \\\hline
两层流体 & 写两层线性式 & 壁面无滑移；界面速度连续；剪应力连续 & 相等的是$\tau$不是$du/dy$ \\\hline
壁律题 & 写$u_\tau,y^+,u^+$ & 由$\tau_w$求$u_\tau$；由$y^+$判粘性底层/对数层；反求$y$或$u$ & $\bar u$和$u(y)$分清 \\\hline
平板边界层 & 写$Re_x=Ux/\nu$ & 判层/湍；求$\delta,C_f$；阻力$D=\bar C_fqA$ & 单面/双面面积 \\\hline
有临界$Re$ & 写$x_{cr}=Re_{cr}\nu/U$ & $L<x_{cr}$全层；否则分段算阻力 & 不要全段套一种 \\\hline
给速度剖面 & 令$\eta=y/\delta$ & 算$\delta^*,\theta,\tau_w$；用动量积分求$\delta(x)$ & 积分别漏$\delta$ \\\hline
外绕球/柱 & 写$Re=UL/\nu$ & 查$C_D$；$D=C_DqA_{\rm front}$ & 面积是迎风面积 \\\hline
小球低$Re$ & 写Stokes阻力 & 终速：重力-浮力=$3\pi\mu dU$；解$U$ & 反查$Re<1$ \\\hline
分离/尾流 & 写逆压梯度 & $dp/dx>0$使近壁减速；$\tau_w=0$为分离点 & 分离导致压差阻力 \\\hline
\end{tabular}
\ftitle{P4 粘性公式链速用表}
\begin{tabular}{|L{0.19\textwidth}|L{0.23\textwidth}|L{0.34\textwidth}|L{0.19\textwidth}|}\hline
已知壁面剪应力 & 写$u_\tau=\sqrt{\tau_w/\rho}$ & 再算$y^+$或摩擦系数$C_f=\tau_w/(0.5\rho U^2)$ & $\tau_w$单位Pa \\\hline
已知压降求管流 & 写$\Delta p/L$ & 层流用H--P直接；湍流用$\Delta p=\lambda L\rho V^2/(2d)$迭代 & $\Delta p$为正损失 \\\hline
已知流量求压降 & 写$V=Q/A$ & $Re\to\lambda\to h_f\to\Delta p=\rho gh_f$ & 管径用内径 \\\hline
边界层厚度 & 写定义$u(\delta)=0.99U$ & 层流平板估$5x/\sqrt{Re_x}$；湍流用经验式 & $\delta$随$x$增大 \\\hline
位移厚度 & 写$\delta^*$定义 & 表示外流被边界层排挤的等效厚度 & 单位长度 \\\hline
动量厚度 & 写$\theta$定义 & 用于动量积分和阻力；$D'=\rho U^2\theta(L)$量级 & 别和角度混 \\\hline
湍流剪应力 & 写$\tau=\mu d\bar u/dy-\rho\overline{u'v'}$ & 近壁粘性项重要；外层雷诺应力重要 & 负号来自脉动相关 \\\hline
粗糙影响 & 写$\varepsilon^+=\varepsilon u_\tau/\nu$ & 小为水力光滑，大为粗糙控制 & 不是所有湍流都粗糙 \\\hline
阻力系数换算 & 写$D=C_DqA$ & $q=0.5\rho U^2$；平板摩擦用$C_f$，钝体用$C_D$ & 系数定义看面积 \\\hline
量级估算 & 写粘性力/惯性力比$1/Re$ & $Re$大粘性只在边界层重要；$Re$小全场粘性重要 & 解释题可用 \\\hline
\end{tabular}
\ftitle{P4 考场纠错表}
\begin{tabular}{|L{0.22\textwidth}|L{0.25\textwidth}|L{0.33\textwidth}|L{0.15\textwidth}|}\hline
层流湍流混 & 第一行算$Re$ & 圆管用$Vd/\nu$，平板用$Ux/\nu$或$UL/\nu$ & 判后选式 \\\hline
平均/最大混 & 看上横线 & 圆管层流$u_{\max}=2\bar u$；湍流另用经验式 & 不直接等同 \\\hline
壁律变量混 & 写三定义 & $u_\tau=\sqrt{\tau_w/\rho}$，$u^+=u/u_\tau$，$y^+=yu_\tau/\nu$ & 无量纲 \\\hline
边界层面积错 & 先定单/双面 & 摩擦阻力用湿面积；外绕压差阻力用迎风面积 & 面积写明 \\\hline
分离原因错 & 写逆压梯度 & $dp/dx>0$导致近壁速度亏损并回流 & 不是因粘性大 \\\hline
Stokes乱用 & 先验$Re<1$ & 若不满足，用$C_D$图或经验关系 & 反查必须写 \\\hline
剪应力方向 & 对相对运动反向 & 壁面对流体阻碍运动，流体对壁面取反 & 画箭头 \\\hline
\end{tabular}
\ftitle{P4 题图触发器：粘性题先判内流/外流，再判速度定义}
\begin{tabular}{|L{0.16\textwidth}|L{0.26\textwidth}|L{0.33\textwidth}|L{0.20\textwidth}|}\hline
\begin{tikzpicture}[scale=.42,baseline=-1ex]\draw (0,0)--(2,0);\draw (0,.8)--(2,.8);\draw[thick] (.2,.1)..controls(1,.65)..(1.8,.1);\draw[->] (.4,.4)--(1.3,.4);\node at(1,-.25){\pictitle{圆管速度分布}};\end{tikzpicture}
& 写“先由$\bar u$算$Re=\bar u d/\nu$判流态” & 层流：$u_{max}=2\bar u,\lambda=64/Re$；湍流：查Moody或用题给壁律 & 题给平均速度不要当最大速度\\\hline
\begin{tikzpicture}[scale=.42,baseline=-1ex]\draw[thick] (0,0)--(2,0);\draw[thick] (0,.75)--(2,.75);\draw[->,red] (.3,.95)--(1.4,.95);\draw[->] (.35,.1)--(.9,.65);\node at(1,-.25){\pictitle{Couette平板}};\end{tikzpicture}
& 写“设$u=u(y)$，由$\mu d^2u/dy^2=dp/dx$” & 无压差则线性；有重力/压差则二次；用两壁无滑移和中间条件定常数 & 剪应力$\tau=\mu du/dy$，看壁面法向\\\hline
\begin{tikzpicture}[scale=.42,baseline=-1ex]\draw (0,0)--(2,0);\draw[->] (-.1,.35)--(.7,.35);\draw (0,0)..controls(.7,.1)and(1.2,.45)..(2,.75);\draw[dashed] (1.5,0)--(1.5,.55);\node at(1,-.25){\pictitle{平板边界层}};\end{tikzpicture}
& 写“外流速度$U_\infty$，局部$Re_x=U_\infty x/\nu$” & 判层/湍/混合；求$\delta$、$C_f$、阻力$D=C_f(1/2\rho U^2)A$ & 面积用湿面积，局部/平均系数别混\\\hline
\begin{tikzpicture}[scale=.42,baseline=-1ex]\draw[->] (0,.35)--(.8,.35);\draw (.95,.35) circle(.35);\draw[->] (1.35,.35)--(1.9,.35);\node at(.95,-.25){\pictitle{球/圆柱阻力}};\end{tikzpicture}
& 写“算$Re=Ud/\nu$，查/选$C_D$” & 阻力$D=C_D(1/2\rho U^2)A$；球/圆柱迎风面积不同，Stokes小$Re$可用$F=3\pi\mu dU$ & $A$ 是投影迎风面积\\\hline
\begin{tikzpicture}[scale=.42,baseline=-1ex]\draw[thick] (0,0)--(2,0);\draw[thick] (0,.8)--(2,.8);\draw (.2,.38)--(1.8,.38);\draw[->,red] (.4,.95)--(1.4,.95);\node at(1,-.25){\pictitle{两层液体}};\end{tikzpicture}
& 写“两层各自线性，界面速度和剪应力连续” & $\tau$ 两层相同；总速度差$U=\tau h_1/\mu_1+\tau h_2/\mu_2$，反求$\tau$ & 厚度和粘度不能交叉写错\\\hline
\end{tabular}
\ftitle{P4 零基础照抄骨架}
\begin{tabular}{|L{0.16\textwidth}|L{0.25\textwidth}|L{0.36\textwidth}|L{0.18\textwidth}|}\hline
圆管层流 & “由$\bar u$算$Re$，若$Re<2300$用H--P” & $u_{max}=2\bar u$，$\lambda=64/Re$，$\Delta p=\lambda(L/d)\rho\bar u^2/2$ & $\bar u$不是$u_{max}$\\\hline
湍流壁律 & “先算$u_*=\sqrt{\tau_w/\rho}$，再用$u^+,y^+$” & $u^+=\bar u/u_*$或$u/u_*$；$y^+=yu_*/\nu$；由题给经验式反求$\tau_w$或厚度 & 经验式变量定义照题\\\hline
平板边界层 & “先算$Re_x=U_\infty x/\nu$” & 层流$\delta=5x/\sqrt{Re_x}$；湍流/混合按题给式；阻力用平均$C_f$ & 局部$x$与全长$L$分清\\\hline
外绕阻力 & “算$Re$，查/选$C_D$” & $D=C_D(1/2\rho U^2)A$；小球$Re<1$可用Stokes & 迎风投影面积\\\hline
剪应力题 & “壁面剪应力$\tau_w=\mu(du/dy)_w$” & 先写速度分布，再求导，最后代壁面位置；方向与速度梯度相关 & 单位Pa\\\hline
\end{tabular}
\newpage

% ================= PAGE 5 =================
\pg{P5 势流叠加、圆柱、镜像、升力}
\begin{multicols}{3}
\blk{1. 势流总判断}
{题写理想、不可压、无旋、势函数、流函数、叠加。}
{首写 $\nabla\times\bm V=0\Rightarrow\bm V=\nabla\phi$；不可压 $\nabla^2\phi=0$。二维不可压：$u=\psi_y,v=-\psi_x$。}
{复杂流场直接叠加：$\phi=\sum\phi_i,\ \psi=\sum\psi_i$；速度由偏导求。}
{Bernoulli 可在无旋势流全场任意两点用；真实粘性阻力不能用势流给。}

\blk{2. 基元流识别表}
{看见源、汇、点涡、偶极子、均匀流。}
{首写对应$\phi,\psi$。}
{均匀流：$\phi=Ur\cos\theta,\psi=Ur\sin\theta$；源：$v_r=Q/(2\pi r),\psi=Q\theta/(2\pi)$；涡：$v_\theta=\Gamma/(2\pi r)$；偶极：$\phi=M\cos\theta/r$。}
{半平面/壁面问题常把$2\pi$换成$\pi$，先看是否有镜像。}

\blk{3. 源+均匀流半体}
{均匀流中有点源，求驻点/分界线/半体形状。}
{首写总流函数 $\psi=Ur\sin\theta+Q\theta/(2\pi)$。}
{驻点：$v_r=0,v_\theta=0$，得 $\theta=\pi,\ r=Q/(2\pi U)$；分界流线过驻点，$\psi=Q/2$。}
{远处半体宽度 $b=Q/U$；方向由源/汇号判。}

\blk{4. 圆柱绕流无环量}
{均匀流绕圆柱、圆柱压力、驻点、达朗贝尔佯谬。}
{首写“均匀流+偶极子叠加，圆柱面$r=a$为$\psi=0$流线”。}
{$\phi=U(r+a^2/r)\cos\theta$，$\psi=U(r-a^2/r)\sin\theta$；$v_r=U(1-a^2/r^2)\cos\theta$；$v_\theta=-U(1+a^2/r^2)\sin\theta$。}
{壁面$C_p=1-4\sin^2\theta$；压力积分阻力升力均为0，只适合理想无粘。}

\blk{5. 圆柱有环量/升力}
{圆柱叠加点涡、机翼升力、Kutta-Joukowski。}
{首写壁面切向速度 $v_\theta=-2U\sin\theta+\Gamma/(2\pi a)$。}
{用 Bernoulli 得压力分布；升力 $L'=\rho U\Gamma$。驻点由 $v_\theta=0$，即 $\sin\theta=\Gamma/(4\pi Ua)$。}
{升力方向按来流和环量右手关系；阻力仍为0是理想结论。}

\blk{6. 镜像法}
{点源/点涡靠近直壁、半平面、壁面吸水。}
{首写“固壁是流线，法向速度为0”。}
{源/汇靠直壁用同号镜像；点涡靠直壁用反号镜像；写总$\psi$后令壁面$\psi=C$，驻点由速度为0。}
{半平面源汇常见分母从$2\pi$变$\pi$；符号错会导致壁面有穿透速度。}

\blk{7. 边角绕流}
{楔角、角域、问速度在角点是否无穷。}
{首写 $\phi=Ar^n\cos n\theta,\psi=Ar^n\sin n\theta$。}
{壁面为$\psi=0$，夹角$\alpha$满足 $n=\pi/\alpha$；速度 $V\sim r^{n-1}$。}
{$\alpha<\pi$速度趋0；$\alpha>\pi$速度趋无穷，实际要靠粘性修正。}

\blk{8. 非定常 Bernoulli}
{无旋非定常、U管振荡、势函数含$t$。}
{首写 $\partial\phi/\partial t+V^2/2+p/\rho+\Pi=f(t)$。}
{同一时刻两点相减消去$f(t)$；若定常则$\partial\phi/\partial t=0$。}
{只适用于无旋；有损失/粘性要改机械能方程。}

\blk{9. Kelvin/Lagrange/Helmholtz}
{证明环量守恒、无旋保持、涡线涡管。}
{首写条件：理想、正压、保守体力。}
{Kelvin：$D\Gamma/Dt=0$；Lagrange：若初始$\Omega=0$，以后仍无旋；Helmholtz：涡线随体、涡管强度守恒、涡线不能在流体内开始/终止。}
{必须写条件；不满足粘性/非正压时不能直接用。}

\blk{10. 极坐标速度关系}
{势流题给$\phi,\psi$含$r,\theta$。}
{首写 $v_r=\phi_r,\ v_\theta=(1/r)\phi_\theta$。}
{流函数：$v_r=(1/r)\psi_\theta,\ v_\theta=-\psi_r$；速度大小$V^2=v_r^2+v_\theta^2$。}
{角向导数必须除$r$；这是最常见失分点。}

\blk{11. 驻点通法}
{求驻点、分界线、物体边界。}
{首写驻点条件 $v_r=0,\ v_\theta=0$ 或 $u=v=0$。}
{先由速度方程求坐标；再把驻点代入$\psi$得到分界流线常数。}
{驻点不是压力最大点的同义词，但理想定常中常对应速度零。}

\blk{12. 压力由速度求}
{势流题问压强、压力系数、升力。}
{首写定常 Bernoulli：$p=p_\infty+\frac12\rho(U_\infty^2-V^2)$。}
{压力系数 $C_p=(p-p_\infty)/(0.5\rho U_\infty^2)=1-(V/U_\infty)^2$。}
{必须先求壁面速度；不要直接对势函数积分压力。}

\blk{13. 压力积分求力}
{圆柱/物体表面压力分布已知，求升阻。}
{首写 $d\bm F=-p\bm n\,dA$。}
{二维单位长度圆柱：$dA=a\,d\theta$；分解到$x,y$后积分$0\to2\pi$。}
{题问流体对物体，用压力指向物体内法线；方向符号要随定义说明。}

\blk{14. 源汇对/Rankine 卵形}
{均匀流+源+汇，求封闭流线、卵形长度宽度。}
{首写总流函数 $\psi=U r\sin\theta+\frac{Q}{2\pi}(\theta_1-\theta_2)$。}
{驻点由总速度为0求；过驻点的$\psi$为物体边界。}
{源汇位置和强度要带符号；别把两者都当源。}

\blk{15. 点涡自由面/压强}
{涡旋、水面下凹、问压强高低或液面形状。}
{首写 $v_\theta=\Gamma/(2\pi r)$。}
{自由面$p=p_a$，用 Bernoulli：$z+\frac{v^2}{2g}=C$，得靠近中心液面降低。}
{中心奇点理论速度无穷，实际有涡核，不能把$r=0$代入。}

\blk{16. 圆柱推导模板}
{题要求由 Laplace 方程推圆柱绕流。}
{首写设 $\phi=F(r)\cos\theta$。}
{代入极坐标 Laplace 得 $r^2F''+rF'-F=0$，解$F=C_1r+C_2/r$；用远场$C_1=U$，壁面$v_r(a)=0$得$C_2=Ua^2$。}
{推导最后写$r=a$是流线/不穿透条件。}

\blk{17. 机翼/库塔条件}
{机翼模型、后缘、环量确定、升力。}
{首写“实际环量由后缘速度有限的 Kutta 条件确定”。}
{有环量后用 $L'=\rho U\Gamma$；若题只问原理，写上表面速度大压强小、下表面反之。}
{没有环量的理想对称绕流不能产生升力。}

\blk{18. 叠加题落笔}
{多个基元流组合，看起来很乱。}
{首写总势/总流函数等于各部分相加。}
{先写$\phi$或$\psi$，再求速度，再用边界/驻点/压力。不要先画流线猜答案。}
{线性叠加只对势函数/流函数/速度成立；压力不能直接叠加。}

\end{multicols}
\ftitle{P5 势流叠加：先写总$\phi/\psi$，再求速度，再用边界/压力}
\begin{tabular}{|L{0.19\textwidth}|L{0.23\textwidth}|L{0.34\textwidth}|L{0.19\textwidth}|}\hline
\textbf{看到/图形} & \textbf{第一行} & \textbf{中间顺序} & \textbf{最后检查}\\\hline
基元流组合 & 写$\phi=\sum\phi_i,\psi=\sum\psi_i$ & 先列每个基元；再求速度；最后找驻点/边界/压力 & 压力不可直接叠加 \\\hline
点源/点汇 & 写$v_r=\pm Q/(2\pi r)$ & 流函数含$\theta$；半平面看是否改$\pi$ & 源正汇负 \\\hline
点涡 & 写$v_\theta=\Gamma/(2\pi r)$ & 速度切向；用Bernoulli求压力；有壁面用镜像 & 原点奇异不代入 \\\hline
均匀流+源 & 写总$\psi$ & 驻点由速度零；过驻点流线为分界线；远处宽度由流量守恒 & 分界线是物体边界 \\\hline
壁面镜像 & 写壁面为流线 & 源汇同号镜像；点涡反号镜像；叠加后检验法向速度0 & 镜像符号最关键 \\\hline
圆柱无环量 & 写均匀流+偶极 & 得$v_\theta(a)=-2U\sin\theta$；$C_p=1-4\sin^2\theta$；积分力 & 阻力0是理想结论 \\\hline
圆柱有环量 & 写$v_\theta(a)=-2U\sin\theta+\Gamma/(2\pi a)$ & 驻点由$v_\theta=0$；升力$L'=\rho U\Gamma$ & 方向按环量判断 \\\hline
求压力分布 & 写$C_p=1-(V/U)^2$ & 先算壁面速度大小；再代Bernoulli；需要力则积分 & 速度大压强小 \\\hline
角域流 & 写$\phi=Ar^n\cos n\theta$ & 壁面$\psi=0$；$n=\pi/\alpha$；速度$\sim r^{n-1}$ & 凹角可能奇异 \\\hline
非定常势流 & 写非定常Bernoulli & 两点相减消$f(t)$；含$\partial\phi/\partial t$ & 只适合无旋 \\\hline
环量定理 & 写理想正压保守体力 & Kelvin守恒；Lagrange无旋保持；Helmholtz涡线性质 & 条件必须写 \\\hline
机翼升力 & 写Kutta条件定环量 & 后缘速度有限$\to\Gamma$；再用$L'=\rho U\Gamma$ & 无环量无升力 \\\hline
\end{tabular}
\ftitle{P5 基元流与常用结果速用表}
\begin{tabular}{|L{0.19\textwidth}|L{0.23\textwidth}|L{0.34\textwidth}|L{0.19\textwidth}|}\hline
均匀流沿$x$ & 写$\phi=Ux,\psi=Uy$ & 极坐标为$\phi=Ur\cos\theta,\psi=Ur\sin\theta$ & 方向变则换角度 \\\hline
点源强度$Q$ & 写$\psi=Q\theta/(2\pi)$ & 两射线间流量$Q\Delta\theta/(2\pi)$ & 全平面才$2\pi$ \\\hline
点汇 & 写源的负号 & 速度向内，$\phi,\psi$均取负源式 & 符号决定驻点位置 \\\hline
点涡环量$\Gamma$ & 写$v_\theta=\Gamma/(2\pi r)$ & 环量$\oint V\cdot dl=\Gamma$ & 原点外无旋 \\\hline
偶极子 & 写$\phi=M\cos\theta/r$ & 常与均匀流叠加成圆柱；$M=Ua^2$ & 方向看$\cos\theta$轴 \\\hline
圆柱驻点 & 写$v_\theta(a)=0$ & 无环量在$\theta=0,\pi$；有环量按$\sin\theta=\Gamma/(4\pi Ua)$ & 驻点可能移位 \\\hline
圆柱压力最大 & 写$V=0$处$p$最大 & 无环量前后驻点$C_p=1$；侧面$C_p=-3$ & 只对理想势流 \\\hline
升力方向 & 写$\bm L'=\rho \bm U\times \bm\Gamma$类比 & 平面题按速度叠加判断哪侧压强低 & 写明方向 \\\hline
边界为流线 & 写$\psi=C$ & 固壁不穿透；自由流线也可作为物体外形 & 常数不重要 \\\hline
速度奇异 & 写理论奇点 & 源/涡/凹角处速度可能无穷，实际由粘性/有限尺寸修正 & 概念题要说明 \\\hline
\end{tabular}
\ftitle{P5 考场纠错表}
\begin{tabular}{|L{0.22\textwidth}|L{0.25\textwidth}|L{0.33\textwidth}|L{0.15\textwidth}|}\hline
把压力叠加 & 立刻停 & 只能叠加$\phi,\psi,\bm V$；压力由总速度进Bernoulli & 先总速度 \\\hline
极坐标少$1/r$ & 写弧长$r d\theta$ & 角向速度公式必须带$1/r$ & 回代检查 \\\hline
镜像符号错 & 看壁面不穿透 & 源汇同号，点涡反号；再验证壁面法向速度0 & 验边界 \\\hline
圆柱阻力不为0 & 看是否理想无粘 & 势流压力积分阻力0；真实阻力来自分离 & 概念分清 \\\hline
升力方向乱 & 看哪侧速度大 & 速度大压强小；环量改变上下速度 & 写方向 \\\hline
驻点找错 & 令速度为零 & 不是令$\phi$或$\psi$为零；分界线才用$\psi=C$ & 两步分开 \\\hline
伯努利乱用 & 看是否无旋 & 势流全场可用；有旋只沿同一流线 & 条件写明 \\\hline
\end{tabular}
\ftitle{P5 题图触发器：势流题看边界，先写叠加式}
\begin{tabular}{|L{0.16\textwidth}|L{0.26\textwidth}|L{0.33\textwidth}|L{0.20\textwidth}|}\hline
\begin{tikzpicture}[scale=.42,baseline=-1ex]\draw[->] (0,.35)--(.75,.35);\node at(.95,.35){$\bullet$};\draw (1.05,.2)..controls(1.45,.2)and(1.7,.5)..(1.95,.75);\node at(1,-.25){\pictitle{均流+点源}};\end{tikzpicture}
& 写“$\psi=Ur\sin\theta+Q\theta/(2\pi)$” & 驻点令$v_r=v_\theta=0$；边界是$\psi=$常数；半体宽度由流量平衡得 & 半平面/全平面分母不同\\\hline
\begin{tikzpicture}[scale=.42,baseline=-1ex]\draw[->] (0,.35)--(.7,.35);\draw (.95,.35) circle(.35);\draw[->] (1.4,.35)--(1.9,.35);\node at(.95,-.25){\pictitle{圆柱无环量}};\end{tikzpicture}
& 写“均流+偶极，圆柱$r=a$为流线” & $v_r=U(1-a^2/r^2)\cos\theta$，表面$v_\theta=-2U\sin\theta$，再用Bernoulli求$p$ & 前后压力对称，阻力为0\\\hline
\begin{tikzpicture}[scale=.42,baseline=-1ex]\draw[->] (0,.35)--(.7,.35);\draw (.95,.35) circle(.35);\draw[->,red] (1.3,.75) arc(55:-230:.36);\node at(.95,-.25){\pictitle{圆柱有环量}};\end{tikzpicture}
& 写“在圆柱绕流上叠加点涡$\Gamma$” & 表面$v_\theta=-2U\sin\theta+\Gamma/(2\pi a)$；上下速度差给升力$L'=\rho U\Gamma$ & 升力方向由环量方向判\\\hline
\begin{tikzpicture}[scale=.42,baseline=-1ex]\draw[thick] (0,0)--(2,0);\node at(.7,.45){$\bullet$};\node at(.7,-.45){$\circ$};\draw[dashed] (.7,-.05)--(.7,-.35);\node at(1,-.75){\pictitle{壁面镜像}};\end{tikzpicture}
& 写“壁面用镜像满足不穿透” & 固壁：同号源/反号涡按边界选；自由面常压条件另判；写总$\phi$或$\psi$再求导 & 镜像不是多一个真实物体\\\hline
\begin{tikzpicture}[scale=.42,baseline=-1ex]\draw (0,0)--(1,.55)--(2,0);\draw[->] (.15,.65)--(.95,.65);\node at(1,-.25){\pictitle{边角绕流}};\end{tikzpicture}
& 写“极坐标设$\phi=Ar^n\cos n\theta,\psi=Ar^n\sin n\theta$” & 用两条固壁流线确定$n$，再由远场/边界确定$A$；速度由求导得 & 角度必须用弧度参与$n$\\\hline
\end{tabular}
\ftitle{P5 零基础照抄骨架}
\begin{tabular}{|L{0.16\textwidth}|L{0.25\textwidth}|L{0.36\textwidth}|L{0.18\textwidth}|}\hline
叠加题 & “由于方程线性，总势/流函数可叠加” & 写$\phi=\sum\phi_i$或$\psi=\sum\psi_i$；速度由求导；边界用$\psi=C$或法向速度0 & 先统一坐标原点\\\hline
驻点/分界线 & “驻点满足$v_x=v_y=0$或$v_r=v_\theta=0$” & 解驻点；取过驻点的流函数值作为分界流线；画边界 & 驻点不一定在原点\\\hline
圆柱压力 & “表面$r=a$，用表面切向速度” & 先得$V(\theta)$；$C_p=1-(V/U)^2$；再积分压力求力/升力 & 不要用径向速度\\\hline
镜像法 & “为满足壁面不穿透，引入镜像奇点” & 写真实+镜像总函数；在壁面检查法向速度为0；再求目标量 & 镜像只为边界服务\\\hline
伯努利求压 & “理想定常不可压：$p=p_\infty+\frac12\rho(U_\infty^2-V^2)$” & 先求局部速度大小，再求压强差；力需沿边界积分 & 高速低压\\\hline
\end{tabular}
\newpage

% ================= PAGE 6 =================
\pg{P6 可压缩、喷管、激波、膨胀波、相似}
\begin{multicols}{3}
\blk{1. 可压缩入门}
{空气、声速、马赫数、压缩性、$Ma$。}
{首写 $Ma=V/a,\ a=\sqrt{\gamma RT}$。}
{$Ma\le0.3$ 常可忽略压缩；否则用可压关系。等熵：$T_0/T=1+(\gamma-1)M^2/2$，$p_0/p=(T_0/T)^{\gamma/(\gamma-1)}$。}
{气体状态方程 $p=\rho RT$ 必须用绝压和绝温 K。}

\blk{2. Laval 喷管判型}
{收缩/缩扩喷管、喉部、背压、阻塞、最大流量。}
{首写比较背压与临界压比。空气临界 $p^*/p_0=0.528,\ T^*/T_0=0.833$。}
{若阻塞：喉部$M=1,A_t=A^*$，$\dot m_{\max}=0.0404A^*p_0/\sqrt{T_0}$。未阻塞：出口内压常等于背压，用等熵压比求$M$。}
{$p_b$是外部背压，不一定等于出口内压$p_e$；阻塞后降背压不增大质量流量。}

\blk{3. 面积--马赫数}
{给 $A/A^*$ 求 $M$ 或给 $M$ 求面积。}
{首写面积关系，并标明有亚声/超声两支。}
{$A/A^*=1$ 时 $M=1$；收缩段亚声加速，扩张段若已超声继续加速。}
{支路由题意决定：喷管入口通常亚声，喉后缩扩喷管可能超声。}

\blk{4. 正激波}
{题写正激波、管内激波、波前波后 1$\to$2。}
{首写“波前等熵求$M_1,p_1,T_1$，过正激波求$M_2,p_2,T_2,p_{02}$”。}
{空气常用：$M_2^2=(1+0.2M_1^2)/(1.4M_1^2-0.2)$；$p_2/p_1=1+\frac{2\gamma}{\gamma+1}(M_1^2-1)$；再用波后等熵。}
{正激波后 $T_0$不变、$p_0$下降；不能继续用波前总压。}

\blk{5. 斜激波}
{超声速遇压缩角、楔形、凹角、斜激波角$\beta$。}
{首写 $M_{n1}=M_1\sin\beta$，法向按正激波。}
{查/解 $\theta$-$\beta$-$M$；弱解常见，强解损失大。波后切向速度不变，法向按正激波变。}
{$\beta=\mu=\arcsin(1/M_1)$ 为马赫波；$\beta=90^\circ$ 为正激波；超过最大折转角会脱体。}

\blk{6. Prandtl--Meyer 膨胀波}
{超声速遇凸角、喷管外膨胀、低压区。}
{首写 $\nu(M_2)=\nu(M_1)+\delta$。}
{查 PM 函数得$M_2$；再用等熵关系求$p_2,T_2,\rho_2$。膨胀波使速度升高、压强温度降低。}
{PM 是连续等熵膨胀，$p_0$不变；不要用激波关系。}

\blk{7. 管外膨胀/压缩波系}
{喷管出口压力和背压不匹配，图有外部波。}
{先由$A_e/A_t$等熵求喷管出口内压$p_e$。}
{若$p_e>p_b$：出口后继续膨胀，PM；若$p_e<p_b$：外部压缩，斜激波/正激波；若$p_e=p_b$：完全膨胀。}
{只在匹配或全亚声特殊情况才可令$p_e=p_b$。}

\blk{8. Fanno 等截面摩擦管}
{可压气体在等截面长管，有摩擦、问最大长度/临界。}
{首写：等截面、绝热、有摩擦，查 Fanno。}
{由入口$M_i$查 $4fL^*/D$；实际 $L$ 与 $L^*$ 比较。摩擦使流动趋向 $M=1$，$T_0$不变，$p_0$下降。}
{这不是等熵；不能用喷管面积关系替代。}

\blk{9. 波阻/声障/声爆概念}
{题问波阻、声障、声爆、激波熵增。}
{首句：超声速流经激波发生不可逆压缩，熵增导致总压损失。}
{波阻来自激波造成的动量损失和压力分布变化；正激波损失最大。声爆是激波传播到地面形成突跃压力。}
{概念题要写“与摩擦不同，属于压缩波不可逆损失”。}

\blk{10. 量纲分析}
{题给多个物理量，要求建立无量纲关系。}
{首写 Buckingham：$n$个量、$r$个基本量纲，得$n-r$个$\pi$。}
{选重复变量：量纲独立、不含待求量、覆盖$M,L,T$；设 $\pi=X\rho^aV^bL^c$ 配平指数。}
{每个$\pi$必须无量纲；重复变量不能全来自同一量纲。}

\blk{11. 模型相似}
{模型实验、缩尺、相似准则、求原型量。}
{先写几何相似，再选主导准则。}
{$Re=\rho VL/\mu$ 主粘性；$Fr=V/\sqrt{gL}$ 主重力；$Ma$ 主压缩；$We$ 主表面张力；$Eu=\Delta p/(\rho V^2)$ 主压强。}
{同流体缩尺常 $Re$ 与 $Fr$ 冲突；保主导准则叫部分相似。}

\blk{12. 水击/水波/明渠}
{题问水击波速、深水波、扰动能否上游传播。}
{水击首写 $\Delta p=\rho a\Delta V$；深水波首写 $h>\lambda/2$ 或 $kh>\pi$；明渠首写 $Fr=V/\sqrt{gh}$。}
{水击 $a\approx\sqrt{K/\rho}/\sqrt{1+KD/(Ee)}$；深水 $c=\sqrt{g\lambda/(2\pi)}$；浅水 $c=\sqrt{gh}$。}
{$Fr<1$缓流扰动可上游，$Fr>1$急流不可上游。}

\blk{13. 等熵任意两点}
{无激波、绝热无摩擦、给两点$p,T,M,V$。}
{首写总温总压不变：$T_{01}=T_{02},p_{01}=p_{02}$。}
{由一个点$M$求$T_0,p_0$；到另一点用面积/压力/速度条件反求$M$，再求$T,p,\rho,V$。}
{遇激波、摩擦、热交换时$p_0$不再守恒。}

\blk{14. 质量流量公式链}
{喷管或气流题求$\dot m$。}
{首写 $\dot m=\rho VA$。}
{等熵代换：$\dot m=A p_0\sqrt{\gamma/(RT_0)}\,M(1+\frac{\gamma-1}{2}M^2)^{-(\gamma+1)/(2(\gamma-1))}$；喉部$M=1$得最大值。}
{用哪个截面都可以，但$\rho,V,A$必须是同一截面状态。}

\blk{15. 背压状态判定}
{缩扩喷管给$p_0,A_e/A_t,p_b$。}
{首写“几何给出候选等熵出口压$p_{e,s}$和临界压”。}
{背压高：全亚声；再降：喉部阻塞；继续降：激波在扩张段并向出口移；激波到出口后为外部压缩波；设计背压$p_b=p_{e,s}$；更低为外部膨胀波。}
{阻塞只说明喉部$M=1$，不说明出口一定超声。}

\blk{16. 内激波定位流程}
{题说激波在喷管内或要求判断位置。}
{先假设激波在某截面$A_s$。}
{波前：由$A_s/A_t$取超声支得$M_1$；过正激波得$M_2,p_{02}$；波后：用新的$p_{02}$和亚声支算到出口，使出口内压等于背压。}
{若算出的$A_s$不在喉后扩张段，假设位置错。}

\blk{17. 激波后继续求出口}
{正激波后还要问出口压力/速度。}
{首写“波后总压换成$p_{02}$继续等熵”。}
{用$M_2$求波后$A^*_2$，再由$A_e/A^*_2$取亚声支求$M_e$，最后求$p_e,T_e,V_e$。}
{不能拿波前$A^*$和$p_0$一路算到出口。}

\blk{18. 斜激波解题链}
{楔角压缩、给$M_1,\theta$求波后。}
{首写查/解$\beta$。}
{顺序：$M_1,\theta\to\beta\to M_{n1}\to$正激波关系求$M_{n2},p_2,T_2\to M_2=M_{n2}/\sin(\beta-\theta)$。}
{弱解/强解若未说明，通常取弱解；超过最大偏转角写脱体。}

\blk{19. 膨胀波解题链}
{凸角、转角$\delta$、求波后$M,p,T$。}
{首写$\nu_2=\nu_1+\delta$。}
{顺序：$M_1\to\nu_1\to\nu_2\to M_2\to$等熵压温比。}
{凸角膨胀；凹角压缩。别把二者公式混。}

\blk{20. 相似题落笔}
{问模型和原型速度、力、压差换算。}
{首写选择主导准则并令模型=原型。}
{由相似准则求速度比；再由目标量无量纲系数相等求力/压差/功率比，如$F\sim \rho V^2L^2$，$P\sim \rho V^3L^2$。}
{同一题可能只要求部分相似；说明忽略的次要效应。}

\blk{21. 可压缩答题检查}
{任何 Ma、喷管、激波题收尾。}
{最后写状态：亚声/临界/超声/阻塞/有无激波。}
{检查：$M<1$扩张减速，$M>1$扩张加速；激波后$M$通常降为亚声；膨胀波$p,T$降、$M$升；激波$p,T$升、$p_0$降。}
{若结果违反这些趋势，优先查支路和总压。}

\end{multicols}
\ftitle{P6 可压缩与相似：先判$Ma$和是否有不可逆过程}
\begin{tabular}{|L{0.19\textwidth}|L{0.23\textwidth}|L{0.34\textwidth}|L{0.19\textwidth}|}\hline
\textbf{看到/图形} & \textbf{第一行} & \textbf{中间顺序} & \textbf{最后检查}\\\hline
给速度和温度 & 写$Ma=V/\sqrt{\gamma RT}$ & 判$Ma\le0.3$能否忽略压缩；若不能，用总量关系 & $T$用K \\\hline
等熵压比 & 写$p_0/p=(1+0.2M^2)^{3.5}$ & 已知压比反求$M$；再求$T,\rho,V$ & 空气指数才是3.5 \\\hline
收缩喷管 & 比较$p_b/p_0$与0.528 & 大于临界未阻塞；小于等于临界喉/出口$M=1$，流量最大 & 降背压不再增流量 \\\hline
缩扩喷管 & 写喉部候选$M=1$ & 用$A_e/A_t$求等熵出口；再和背压比判全亚声/内激波/外波/设计 & 出口内压不总等于背压 \\\hline
求最大质量流量 & 写喉部$A^*$ & 用$\dot m^*=0.0404A^*p_0/\sqrt{T_0}$或通式$M=1$ & 只适合空气常数 \\\hline
正激波 & 写波前等熵、波内激波、波后等熵 & $M_1\to M_2,p_2/p_1,T_2/T_1,p_{02}$；再继续下游 & $p_0$降，$T_0$不变 \\\hline
斜激波 & 写$M_{n1}=M_1\sin\beta$ & 解$\beta$；法向按正激波；$M_2=M_{n2}/\sin(\beta-\theta)$ & 通常取弱解 \\\hline
膨胀角 & 写$\nu_2=\nu_1+\delta$ & 查PM得$M_2$；等熵求$p,T$ & 膨胀$p,T$降，$M$升 \\\hline
Fanno管 & 写等截面绝热摩擦 & 查$4fL^*/D$；与实际$L$比；趋向$M=1$ & 不是等熵 \\\hline
量纲分析 & 写$n-r$个$\pi$ & 选重复变量；配指数；写无量纲关系 & 每个$\pi$无量纲 \\\hline
模型相似 & 写几何相似+主准则 & 按$Re/Fr/Ma/We/Eu$求比例；再换算力/压差/功率 & 说明部分相似 \\\hline
水击/明渠 & 写对应判据 & 水击$\Delta p=\rho a\Delta V$；明渠$Fr=V/\sqrt{gh}$；水波看深浅水 & $Fr$判扰动方向 \\\hline
\end{tabular}
\ftitle{P6 空气常用关系与判错表}
\begin{tabular}{|L{0.19\textwidth}|L{0.23\textwidth}|L{0.34\textwidth}|L{0.19\textwidth}|}\hline
空气等熵 & 写$T_0/T=1+0.2M^2$ & $p_0/p=(T_0/T)^{3.5}$，$\rho_0/\rho=(T_0/T)^{2.5}$ & 只对$\gamma=1.4$ \\\hline
临界状态 & 写$M=1$ & $T^*/T_0=0.833$，$p^*/p_0=0.528$，$\rho^*/\rho_0=0.634$ & 不是所有喉压比 \\\hline
声速 & 写$a=\sqrt{\gamma RT}$ & 先把摄氏度转K；$V=Ma$ & $R=287$用于空气 \\\hline
总压损失 & 写$p_{02}<p_{01}$ & 激波/摩擦会降总压；等熵和PM不降总压 & $T_0$可不变 \\\hline
超声扩张段 & 写面积增大加速 & $M>1$时$A\uparrow,V\uparrow,p\downarrow$ & 和亚声相反 \\\hline
亚声扩张段 & 写面积增大减速 & $M<1$时$A\uparrow,V\downarrow,p\uparrow$ & 喷管入口常亚声 \\\hline
正激波趋势 & 写$M_2<1$ & $p,T,\rho$升；$V,M$降；$p_0$降 & 熵增不可逆 \\\hline
膨胀波趋势 & 写等熵加速 & $M,V$升；$p,T,\rho$降；$p_0$不变 & 只在超声凸角 \\\hline
斜激波趋势 & 写压缩折转 & 流向转向壁面；压强升；可能仍超声 & 角太大脱体 \\\hline
相似比例 & 写主导准则相等 & $Re$给速度比；力比用$\rho V^2L^2$；功率比用$\rho V^3L^2$ & 几何比先定 \\\hline
\end{tabular}
\ftitle{P6 考场纠错表}
\begin{tabular}{|L{0.22\textwidth}|L{0.25\textwidth}|L{0.33\textwidth}|L{0.15\textwidth}|}\hline
压力用错 & 全部改绝压 & 可压缩所有$p,p_0,p_b,p_e$均用绝压 & 单位Pa \\\hline
温度用摄氏 & 改为K & $T(K)=t(^\circ C)+273.15$ & 等熵必须K \\\hline
出口等于背压 & 先判断状态 & 全亚声常可$p_e=p_b$；超声不匹配会有外部波 & 不默认相等 \\\hline
阻塞理解错 & 写喉部$M=1$ & 只代表质量流量达最大，不代表全管超声 & 看扩张段 \\\hline
激波后继续等熵 & 换新总压 & 波后用$p_{02}$，不能沿用$p_{01}$ & $p_0$下降 \\\hline
面积比两支错 & 标亚/超声支 & 收缩段入口多亚声；喉后扩张才可能超声 & 由题意选支 \\\hline
斜激波无解 & 查最大偏转角 & 若$\theta>\theta_{\max}$，写脱体激波 & 不强行算 \\\hline
相似准则冲突 & 写部分相似 & 说明保留主导力，其他效应不能完全满足 & 合理说明 \\\hline
\end{tabular}
\ftitle{P6 题图触发器：可压缩题看图直接分支}
\begin{tabular}{|L{0.16\textwidth}|L{0.26\textwidth}|L{0.33\textwidth}|L{0.20\textwidth}|}\hline
\begin{tikzpicture}[scale=.42,baseline=-1ex]\draw (0,.15)..controls(.6,.45)and(1.1,.45)..(1.8,.15);\draw (0,.95)..controls(.6,.65)and(1.1,.65)..(1.8,.95);\draw[dashed] (.9,.2)--(.9,.9);\node at(.9,-.25){\pictitle{Laval喉部}};\end{tikzpicture}
& 写“先看喉部是否$M=1$，再看背压分支” & 用$p_b/p_0$或面积比判断：全亚声、阻塞、内激波、外激波、设计膨胀 & $0.528$只对应空气临界静压比\\\hline
\begin{tikzpicture}[scale=.42,baseline=-1ex]\draw[->] (0,.45)--(.75,.45);\draw[very thick,red] (1,.05)--(1,.85);\draw[->] (1.25,.45)--(1.85,.45);\node at(.95,-.25){\pictitle{正激波}};\end{tikzpicture}
& 写“波前1超声，波后2亚声，$T_0$不变$p_0$降” & 先由上游$M_1$查/算$M_2,p_2/p_1,T_2/T_1$，再接波后等熵 & 不能只靠$Ma_2$随便算全部\\\hline
\begin{tikzpicture}[scale=.42,baseline=-1ex]\draw[->] (0,.45)--(.8,.45);\draw (1,.15)--(1.8,.15);\draw (1,.15)--(1.8,.55);\draw[red] (1.15,.05)--(1.6,.85);\node at(1,-.25){\pictitle{斜激波楔角}};\end{tikzpicture}
& 写“给楔角$\theta$，先求波角$\beta$” & $M_{n1}=M_1\sin\beta$，按正激波算法向，再由$M_2=M_{n2}/\sin(\beta-\theta)$ & 默认弱解；角过大写脱体\\\hline
\begin{tikzpicture}[scale=.42,baseline=-1ex]\draw[->] (0,.45)--(.8,.45);\draw (1,.15)--(1.8,.15);\draw (1,.15)--(1.8,-.25);\foreach \x in {1.15,1.3,1.45,1.6}{\draw[red] (1,.15)--(\x,-.25);} \node at(1,-.55){\pictitle{凸角膨胀扇}};\end{tikzpicture}
& 写“凸角超声膨胀：$\nu_2=\nu_1+\delta$” & 查PM函数得$M_2$，再用等熵关系算$p_2,T_2$ & 膨胀$p,T$降，$M$升\\\hline
\begin{tikzpicture}[scale=.42,baseline=-1ex]\draw (0,.2)--(2,.2);\draw (0,.75)--(2,.75);\draw[->] (.2,.48)--(1.8,.48);\node at(1,-.25){\pictitle{Fanno长管}};\end{tikzpicture}
& 写“等截面绝热有摩擦，不能用等熵” & 查$4fL^*/D$，比较实际$L$；流动趋向$M=1$，可能阻塞 & 摩擦使$p_0$下降\\\hline
\end{tabular}
\ftitle{P6 零基础照抄骨架}
\begin{tabular}{|L{0.16\textwidth}|L{0.25\textwidth}|L{0.36\textwidth}|L{0.18\textwidth}|}\hline
等熵段 & “本段无激波无摩擦，$p_0,T_0$不变” & 由$M$求$T/T_0,p/p_0,\rho/\rho_0$；或由面积比取亚/超声支 & 支路由物理位置决定\\\hline
阻塞判断 & “先判喉部是否达到$M=1$” & 收缩管比$p_b/p_0$；Laval先喉部临界，再看背压状态图 & 阻塞后流量不再随背压增大\\\hline
激波分段 & “波前等熵$\to$正/斜激波$\to$波后等熵” & 波前用$p_{01}$，波后换成$p_{02}$；若在喷管内，出口用波后总压继续算 & 不可跨激波等熵\\\hline
膨胀波 & “凸角超声膨胀，PM函数相加” & $M_1\to\nu_1$，$\nu_2=\nu_1+\delta$，查$M_2$，用等熵求$p_2,T_2$ & 只有超声凸角\\\hline
相似换算 & “确定主导力，令对应无量纲数相等” & 几何比先定；由$Re/Fr/Ma$求速度比；再用$F\sim\rho V^2L^2$等换算 & 写明忽略哪些效应\\\hline
\end{tabular}
\end{document}
